\subsection{分式环}

定理4. 设A是一个交换环，S是A的一个子集。设 \(\mathbf{A}_{\mathrm{S}}\) 为A的分式幺半群（仅配备乘法），其分母在S中（§2，第4号）。设 \(\varepsilon: \mathbf{A} \to \mathbf{A}_{\mathrm{S}}\) 为典范态射。在 \(\mathbf{A}_{\mathrm{S}}\) 上存在一种且仅一种加法，满足以下条件：

\begin{enumerate}
\def\labelenumi{(\alph{enumi})}
\item
  \(A_{s}\) ，连同这个加法和乘法，构成一个交换环；
\item
  \(\varepsilon\) 是一个环同态。
\end{enumerate}

假设已为 \(\mathbf{A}_{\mathrm{s}}\) 找到满足条件(a)和(b)的一个加法。

设 \(x, y \in A_s\) 。设 \(S'\) 是 \(A\) 的由 \(S\) 生成的稳定乘法子幺半群。存在 \(a, b \in A\) 和 \(p, q \in S'\) ，使得 \(x = a/p, y = b/q\) 。则

\[
x = \varepsilon (a q) \varepsilon (p q) ^ {- 1}, \quad y = \varepsilon (b p) \varepsilon (p q) ^ {- 1},
\]

从而

\begin{enumerate}
\def\labelenumi{(\arabic{enumi})}
\setcounter{enumi}{22}
\tightlist
\item
\end{enumerate}

\[
\begin{array}{r l} x + y & = (\varepsilon (a q) + \varepsilon (b p)) \varepsilon (p q) ^ {- 1} \\ & = \varepsilon (a q + b p) \varepsilon (p q) ^ {- 1} \\ & = (a q + b p) / p q. \end{array}
\]

这证明了加法的唯一性。

我们现在在 \(A_{s}\) 上通过设定 \(x + y = (aq + bp)/pq\) 来定义一个加法。有必要证明这个定义不依赖于a、b、p、q的选取。现在，若 \(a', b' \in A\) , \(p', q' \in S'\) 满足 \(x = a'/p', y = b'/q'\) ，则存在属于 \(S'\) 的 s 和 t，使得 \(ap's = a'ps\) , \(bq't = b'qt\) ，因此

\[
(a q + b p) \left(p ^ {\prime} q ^ {\prime}\right) (s t) = \left(a ^ {\prime} q ^ {\prime} + b ^ {\prime} p ^ {\prime}\right) (p q) (s t)
\]

从而

\[
(a q + b p) / p q = (a ^ {\prime} q ^ {\prime} + b ^ {\prime} p ^ {\prime}) / p ^ {\prime} q ^ {\prime}.
\]

容易验证 \(A_{s}\) 中的加法满足结合律和交换律，0/1是加法的单位元， \((-a)/p\) 是a/p的负元，并且 \(x(y + z) = xy + xz\) 对一切 \(x, y, z \in A_{s}\) 成立。若 \(a, b \in A\) ，则

\[
\varepsilon (a + b) = (a + b) / 1 = a / 1 + b / 1 = \varepsilon (a) + \varepsilon (b)
\]

从而 \(\varepsilon\) 是一个环同态。

定义8. 定理4中定义的环称为与S相伴的分式环，或分母在S中的分式环，并记为A{[}S \(^{-1}\) {]}.

\(\mathbf{A}[\mathbf{S}^{-1}]\) 的零元是 \(0 / 1\) ， \(\mathbf{A}[\mathbf{S}^{-1}]\) 的单位元是 \(1 / 1\) .

我们将在《交换代数》第二章，§ 2 中回到 \(\mathbf{A}[\mathrm{S}^{-1}]\) 的性质。

若S是A中可消去元素的集合，则环 \(A[S^{-1}]\) 称为A的全分式环。此时A被等同于 \(A[S^{-1}]\) 的一个子环，所借助的映射 \(\varepsilon\) 此时是单射的（I，§2，第4号，命题6）。

定理5. 设A是一个交换环，S是A的一个子集，B是一个环，f是A到B中的一个同态，使得 \(f(\mathbf{S})\) 的每个元素都可逆。存在一个且仅一个 \(\bar{f}\) ，它把 \(A[S^{-1}]\) 映入 B，使得 \(f = \bar{f} \circ \varepsilon\) .

我们知道（§2，第4号，定理1）存在一个且仅一个态射 \(\bar{f}\) ，它把乘法幺半群 A{[}S \(^{-1}\) {]} 映入乘法幺半群 B，使得 \(f = \bar{f} \circ \varepsilon\) 。设 \(a, b \in \mathbf{A}, p, q \in \mathbf{S}'\) （由S生成的A的稳定乘法子幺半群）。由于 \(f(A)\) 的元素两两交换，

\[
\begin{array}{r l} \bar {f} (a | \dot {p} + b | q) & = \bar {f} ((a q + b \dot {p}) | \dot {p} q) = f (a q + b \dot {p}) f (\dot {p} q) ^ {- 1} \\ & = (f (a) f (q) + f (b) f (\dot {p})) f (\dot {p}) ^ {- 1} f (q) ^ {- 1} \\ & = f (a) f (\dot {p}) ^ {- 1} + f (b) f (q) ^ {- 1} \\ & = \bar {f} (a | \dot {p}) + \bar {f} (b | q). \end{array}
\]

因此 \(\tilde{f}\) 是一个环同态。

\section{域}

\subsection{域}

定义 1. 环 K 称为域，如果它不仅仅由 0 组成，并且 K 的每个非零元素都可逆。

域K的非零元素集合在乘法下构成一个群，它正是环K的乘法群 \(K^{*}\) （ \(\S\) 8，第1号）。一个域的反环仍是域。若一个域的乘法是交换的，则称该域为交换域；这样的域与其反环等同。非交换域有时称为除环（体）。

例子。*(1) 我们将在第4号中定义有理数域；在《一般拓扑学》中将定义实数域（《一般拓扑学》，IV，§1，第3号）、复数域（《一般拓扑学》，VIII，§1，第1号）以及四元数域（《一般拓扑学》，VIII，§1，第4号）。除四元数域外，这些域都是交换的。* (2) 环Z/2Z显然是一个域。

设 K 为一个域。K 的每个本身是域的子环 L 称为 K 的子域，此时 K 称为 L 的扩域；这等价于说 L 是 K 的子环，并且对 L 的每个非零元素 x 有 \(x^{-1} \in L\) 。若 \((\mathbf{L}_{i})_{i \in I}\) 是K的一族子域，则 \(\bigcap_{i \in I} L_{i}\) 是 K 的一个子域；因此，对于 K 的每个子集 X，存在包含 X 的 K 的最小子域；称它为由 X 生成的。

命题1. 设 K 是一个域。对于 K 的每个子集 X，X 的中心化子（§8，第5号，例3）X' 是 K 的一个子域。

我们知道（同上引文） \(\mathbf{X}'\) 是 \(\mathbf{K}\) 的子环。另一方面，若 \(x \neq 0\) 与 \(z \in \mathbf{X}\) 可交换，则 \(x^{-1}\) 亦然（I，§2，第3号，命题5），从而 \(\mathbf{X}'\) 包含 \(\mathbf{X}'\) 中每个非零元素的逆。

推论。域K的中心是K的一个（交换）子域。

定理1. 设A是一个环。以下条件等价：

\begin{enumerate}
\def\labelenumi{(\alph{enumi})}
\item
  A 是一个域；
\item
  A 不约化为 0，且 A 仅有的左理想是 0 和 A。
\end{enumerate}

假设A是一个域。那么A不约化为0。设 a 是 A 的一个异于 0 的左理想。存在一个非零的 a 属于 a。对于所有 \(x \in A\) , \(x = (xa^{-1})a \in a\) ；因此 a = A。

假设 A 满足条件 (b)。设 \(x \neq 0\) 属于 A。我们需要证明 x 是可逆的。左理想 Ax 包含 x，因此不为零，从而 Ax = A。于是存在 \(x' \in A\) 使得 \(x'x = 1\) 。则 \(x' \neq 0\) ，因为 \(1 \neq 0\) 。将上述结果应用于 \(x'\) ，可以看出存在 \(x'' \in A\) 使得 \(x''x' = 1\) 。则 \(x'' = x'' \cdot 1 = x''x'x = 1 \cdot x = x\) ，因此 \(xx' = 1\) 。因此 \(x'\) 是 x 的逆。

注。在定理1中，左理想可以换成右理想。在第八章，§5，第2号中，我们将研究没有异于0和A的双边理想的非零环A；这样的环（称为拟单环）不一定是域 *（例如，以有理数为系数的2阶方阵组成的环 \(\mathbf{M}_2(\mathbf{Q})\) 是拟单的，但不是域）*。

推论1. 设A是一个环，且a是A的一个双边理想。要使环A/a成为域，当且仅当a是A的一个极大左理想。

A/a的左理想具有b/a的形式，其中b是A的包含a的左理想（§8，第8号，命题5的推论）。说A/a ≠ 0意味着a ≠ A。在此假设下，A/a是域当且仅当A的包含a的左理想仅有a和A（定理1），由此得出推论。

推论2. 设A是一个非零交换环。则存在一个从A到某个交换域上的同态。

由 Krull 定理（§ 8，第 6 号，定理 1），在 A 中存在一个极大理想 \(a\) 。则 A/a 是一个域（推论 1）。

推论 3. 设 a 为整数 \(\geqslant 0\) 。环 \(\mathbf{Z} / a\mathbf{Z}\) 为域当且仅当 a 为素数。

这由推论1及 §8 第11号得出。

对于素数p，域Z/pZ记作 \(F_{p}\) .

定理2. 设K是一个域，A是一个非零环。若f是K到A中的一个同态，则A的子环 \(f(\mathbf{K})\) 是一个域，且f定义了K到 \(f(\mathbf{K})\) 上的一个同构。

设 \(\mathfrak{a}\) 为 \(f\) 的核。则 \(1 \notin \mathfrak{a}\) ，因为 \(f(1) = 1 \neq 0\) 在 \(\mathbf{A}\) 中，并且由于 \(\mathfrak{a}\) 是 \(\mathbf{K}\) 的左理想，由定理1有 \(\mathfrak{a} = \{0\}\) 。因此 \(f\) 是单射，从而是 \(\mathbf{K}\) 到子环 \(f(\mathbf{K})\) 上的同构，该子环位于 \(\mathbf{A}\) 中；因此后一环是一个域。

\subsection{整环}

定义2. 环A称为整环（或整区），如果它是交换的、非零的，并且A中两个非零元素的乘积是非零的。

有理整数环Z是整环；它是交换的且非零的；两个整数 \textgreater0 的乘积是非零的；每个非零整数都具有a或-a的形式，其中a \textgreater{} 0 且 \((-a)b = -ab\) , \((-a)(-b) = ab\) 对于 a \textgreater{} 0，b \textgreater{} 0，由此得出我们的断言。

每个交换域都是整环。整环的子环是整环。特别地，交换域的子环是整环。我们将证明反之每个整环 A 都同构于某个交换域的子环。回顾（ \(\S 8\) ，第 12 号）A 已被等同于其全分式环的一个子环。于是我们的断言由以下命题得出：

命题2. 若A是整环，则A的全分式环K是交换环。

环 K 是交换的。它是非零的，因为 \(A \neq \{0\}\) 。由于 A 是整环，A 的每个非零元素都是可消去的，且 K 由形如 a/b 的分数组成，其中 \(b \neq 0\) 。现在 \(a/b \neq 0\) 蕴含 \(a \neq 0\) ，且分数 b/a 即为 a/b 的倒数。

整环的全分式环称为其分式域。这样的环与其在分式域中的像等同。

命题3. 设B为非零环，A为B的一个交换子环，使得A中每个非零元素在B中都可逆。

\begin{enumerate}
\def\labelenumi{(\alph{enumi})}
\item
  A 是一个整环。
\item
  设 \(\mathbf{A}'\) 为 \(\mathbf{A}\) 的分式域。从 \(\mathbf{A}\) 到 \(\mathbf{B}\) 的典范嵌入可以唯一地延拓为一个同构 \(f\) ，它把 \(\mathbf{A}'\) 映到 \(\mathbf{B}\)
\item
  的一个子域上。 \(f(\mathbf{A}')\) 的元素就是形如 \(xy^{-1}\) 的元素，其中 \(x \in \mathbf{A}, y \in \mathbf{A}, y \neq 0\) .
\end{enumerate}

。断言 (a) 是显然的。从 A 到 B 的典范嵌入唯一地延拓为一个同态 \(f\) ，它把 \(\mathbf{A}'\) 映入 \(\mathbf{B}\) （§ 8，第 12 号，定理 5）。断言 (b) 随后由第 1 号定理 2 得出。\(\mathbf{A}'\) 的元素是形如 \(x / y\) 的分数，其中 \(x \in \mathbf{A}, y \in \mathbf{A}, y \neq 0\) ，且 \(f(x / y) = xy^{-1}\) ，由此得 (c)。

\subsection{素理想}

命题4. 设A是一个交换环，p是A的一个理想。以下条件等价：

\begin{enumerate}
\def\labelenumi{(\alph{enumi})}
\item
  环 A/p 是一个整环；
\item
  \(\mathbf{A} \neq \mathfrak{p}\) 并且，如果 \(x \in \mathbf{A} - \mathfrak{p}\) 和 \(y \in \mathbf{A} - \mathfrak{p}\) ，则 \(xy \in \mathbf{A} - \mathfrak{p}\) .
\item
  \(\mathfrak{p}\) 是 \(\mathbf{A}\) 到一个域中的同态的核。
\end{enumerate}

蕴含关系 (c) \(\Rightarrow\) (b) \(\Rightarrow\) (a) 是显然的。若 A/p 是整环，设 \(f\) 为 A/p 到其分式域中的典范嵌入， \(g\) 为 A 到 A/p 上的典范同态；则 p 是 \(f \circ g\) 的核，从而蕴含关系 (a) \(\Rightarrow\) (c) 成立。

定义3. 在交换环A中，满足命题4条件的理想p称为素理想。

示例。(1) 设 A 是一个交换环。若 m 是 A 的一个极大理想，则 m 是素理想；环 A/m 是一个域（第 1 号，定理 1 的推论 1）。

\begin{enumerate}
\def\labelenumi{(\arabic{enumi})}
\setcounter{enumi}{1}
\tightlist
\item
  若 \(\mathbf{A}\) 是整环，则 \(\{0\}\) 是 \(\mathbf{A}\) 的素理想（但一般不是极大理想，环 \(\mathbf{Z}\) 的例子便证明了这一点）。
\end{enumerate}

\subsection{有理数域}

定义4. 整数环Z的分式域称为有理数域，记作Q。Q的元素称为有理数。

因此，每个有理数都具有形式 \(a / b\) ，其中 \(a\) 与 \(b\) 为满足 \(b \neq 0\) 的有理整数（我们甚至可以取 \(b > 0\) ，这由关系

\[
a / b = (- a) / (- b)
\]

所证明）。 \(\mathbf{Q}_{+}\) 用来表示形如 \(a / b\) （其中 \(a\in \mathbf{N}\) ， \(b\in \mathbf{N}^*\) ）的有理数所成的集合。

我们有如下关系：

\begin{enumerate}
\def\labelenumi{(\arabic{enumi})}
\tightlist
\item
\end{enumerate}

\[
\mathbf {Q} _ {+} + \mathbf {Q} _ {+} = \mathbf {Q} _ {+}\tag{2}
\]

\[
\mathbf {Q} _ {+}. \mathbf {Q} _ {+} = \mathbf {Q} _ {+}\tag{3}
\]

\[
\mathbf {Q} _ {+} \cap (- \mathbf {Q} _ {+}) = \{0 \}\tag{4}
\]

\[
\mathbf {Q} _ {+} \cup (- \mathbf {Q} _ {+}) = \mathbf {Q}\tag{5}
\]

\[
\mathbf {Q} _ {+} \cap \mathbf {Z} = \mathbf {N}.
\]

前两个关系由公式 \(a / b + a' / b' = (ab' + a'b) / bb'\) , \((a / b)(a'/b') = aa'/bb'\) , \(0 \in \mathbf{Q}_+\) , \(1 \in \mathbf{Q}_+\) 以及 \(\mathbf{N}\) 对加法和乘法封闭、\(\mathbf{N}^*\) 对乘法封闭这一事实得出。显然 \(0 \in \mathbf{Q}_+\) ，从而 \(0 \in (-\mathbf{Q}_+)\) ；设 \(x\) 属于 \(\mathbf{Q}_+ \cap (-\mathbf{Q}_+)\) ；则存在正整数 \(a, b, a', b'\) 满足 \(b \neq 0\) , \(b' \neq 0\) 且 \(x = a / b = -a'/b'\) ；于是 \(ab' + ba' = 0\) ，从而 \(ab' = 0\) （因为 \(ab' \geqslant 0\) 且 \(ba' \geqslant 0\) ），进而得 \(a = 0\) ，因为 \(b' \neq 0\) ；换言之， \(x = 0\) 。最后，显然 \(\mathbf{N} \subset \mathbf{Z}\) 且 \(\mathbf{N} \subset \mathbf{Q}_+\) 。反之，若 \(x\) 属于 \(\mathbf{Z} \cap \mathbf{Q}_+\) ，则它是一个有理整数；存在两个有理整数 \(a\) 和 \(b\) 满足 \(a \geqslant 0\) , \(b > 0\) 且 \(x = a / b\) ，于是 \(a = bx\) ；若 \(x \notin \mathbf{N}\) ，则 \(-x > 0\) ，于是 \(-a = b(-x) > 0\) ，从而 \(a < 0\) ，与假设矛盾。

给定两个有理数 \(x\) 和 \(y\) ，我们记 \(x \leqslant y\) ，若 \(y - x \in \mathbf{Q}_+\) 。由公式(1)、(3)和(4)容易推出 \(x \leqslant y\) 是

Q上的一个全序，由(5)可知该关系在Z上诱导出通常的序关系。最后，由(1)可得关系 \(x \leqslant y\) 和 \(x' \leqslant y'\) 蕴含 \(x + x' \leqslant y + y'\) ，由(2)可得关系 \(x \leqslant y\) 蕴含 \(xz \leqslant yz\) 对一切 \(z \geqslant 0\) ，且蕴含 \(xz \geqslant yz\) 对 \(z \leqslant 0\) （参见VI，§2，第1号）。

设 \(x\) 是一个有理数。称 \(x\) 为正的，若 \(x \geqslant 0\) ；称其为严格正的，若 \(x > 0\) ；称其为负的，若 \(x \leqslant 0\) ；称其为严格负的，若 \(x < 0\) 。正有理数的集合是 \(\mathbf{Q}_{+}\) ，负有理数的集合是 \(-\mathbf{Q}_{+}\) 。若 \(\mathbf{Q}^{*}\) 表示非零有理数的集合，则严格正数的集合 \(\mathbf{Q}_{+}^{*}\) 等于 \(\mathbf{Q}^{*} \cap \mathbf{Q}_{+}\) ，而 \(-\mathbf{Q}_{+}^{*}\) 是严格负数的集合。

集合 \(\mathbf{Q}_{+}^{*}\) 和 \(\{1, -1\}\) 是乘法群 \(\mathbf{Q}^{*}\) 的子群。每个有理数 \(x \neq 0\) 都能以且仅以一种方式表示成 \(1.y\) 或 \((-1).y\) 的形式，其中 \(y > 0\) ；因此乘法群 \(\mathbf{Q}^{*}\) 是子群 \(\mathbf{Q}_{+}^{*}\) 与 \(\{1, -1\}\) 的乘积， \(x\) 在 \(\mathbf{Q}_{+}^{*}\) 中的分量称为 \(x\) 的绝对值，记作 \(|x|\) ； \(x\) 在 \(\{-1, 1\}\) 中的分量（等于 1 若 \(x > 0\) ，等于 \(-1\) 若 \(x < 0\) ）称为 \(x\) 的符号，记作 \(\operatorname{sgn} x\) .

通常，我们把这两个函数延拓到整个 \(\mathbf{Q}\) 上，办法是设定 \(|0| = 0\) 和 \(\operatorname{sgn} 0 = 0\) 。

\section{逆向极限与正向极限}

在本段中，I 表示一个非空预序集， \(\alpha \leqslant \beta\) 表示 I 上的预序。相对于指标集 I 的集合的逆向（相应地，正向）系统的概念在《集合论》III，§ 7，第 1 号（相应地，《集合论》III，§ 7，第 5 号，在 I 是右有向的假设下）中定义。

\subsection{广群的逆向系统}

定义1. 关于指标集I的广群逆向系统是一个集合的逆向系统 \((\mathbf{E}_{\alpha}, f_{\alpha\beta})\) ，其中每个 \(E_{\alpha}\) 具有一个广群结构，且每个 \(f_{\alpha\beta}\) 都是广群同态。

设 \((\mathbf{E}_{\alpha}, f_{\alpha \beta})\) 是一个其律以乘法记法书写的广群逆向系统。逆极限集 \(\mathbf{E} = \varprojlim \mathbf{E}_{\alpha}\) 是积广群 \(\prod_{\alpha \in I} \mathbf{E}_{\alpha}\) 的子集，由这样的族 \((x_{\alpha})_{\alpha \in I}\) 组成：满足 \(x_{\alpha} = f_{\alpha \beta}(x_{\beta})\) （对 \(\alpha \leqslant \beta\) ）。若 \((x_{\alpha})\) 与 \((y_{\alpha})\) 属于 E，则对 \(\alpha \leqslant \beta\) 有 \(x_{\alpha} = f_{\alpha \beta}(x_{\beta})\) 与 \(y_{\alpha} = f_{\alpha \beta}(y_{\beta})\) ，因此 \(x_{\alpha}y_{\alpha} = f_{\alpha \beta}(x_{\beta})f_{\alpha \beta}(y_{\beta}) = f_{\alpha \beta}(x_{\beta}y_{\beta})\) ；因此 E 是 \(\prod_{\alpha \in I} \mathbf{E}_{\alpha}\) 的一个子广群。将给 E 赋予由 \(\prod_{\alpha \in I} \mathbf{E}_{\alpha}\) 上的律诱导的律；如此得到的广群称为诸广群 \(E_{\alpha}\) 的逆极限广群。它满足以下泛性质：

\begin{enumerate}
\def\labelenumi{(\alph{enumi})}
\item
  对于所有 \(\alpha \in \mathbf{I}\) ，典范映射 \(f_{\alpha}\) （它把 E 映入 \(\mathrm{E}_{\alpha}\) ）是广群同态，且 \(\mathrm{E}_{\alpha} \cdot f_{\alpha} = f_{\alpha \beta} \circ f_{\beta}\) 对 \(\alpha \leqslant \beta\) 成立。
\item
  假设给定一个广群 F 以及同态 \(u_{\alpha} \colon \mathrm{F} \to \mathrm{E}_{\alpha}\) ，满足 \(u_{\alpha} = f_{\alpha \beta} \circ u_{\beta}\) （对 \(\alpha \leqslant \beta\) ）。存在一个且仅一个同态 \(u \colon \mathrm{F} \to \mathrm{E}\) ，使得 \(u_{\alpha} = f_{\alpha} \circ u\) 对所有 \(\alpha \in \mathrm{I}\) 成立（即 \(x \mapsto u(x) = (u_{\alpha}(x))_{\alpha \in \mathrm{I}}\) ）。
\end{enumerate}

如果诸广群 \(E_{\alpha}\) 是结合的（相应地，交换的），那么 E 也是。假设每个广群 \(E_{\alpha}\) 有一个单位元 \(e_{\alpha}\) ，且诸同态 \(f_{\alpha\beta}\) 是保单位的。则 \(e = (e_{\alpha})_{\alpha \in I}\) 属于 E，因为有 \(e_{\alpha} = f_{\alpha\beta}(e_{\beta})\) （对 \(\alpha \leqslant \beta\) ），且它是广群 E 的单位元；用上述记号，诸同态 \(f_{\alpha}\) 是保单位的，且若诸 \(u_{\alpha}\) 保单位，则 u 也保单位。进一步，E 的元素 \(x = (x_{\alpha})_{\alpha \in I}\) 可逆当且仅当每个 \(x_{\alpha}\) 在相应的广群 \(E_{\alpha}\) 中可逆，且 \(x^{-1} = (x_{\alpha}^{-1})_{\alpha \in I}\) ；这由公式 \(f_{\alpha\beta}(x_{\beta}^{-1}) = f_{\alpha\beta}(x_{\beta})^{-1} = x_{\alpha}^{-1}\) （对 \(\alpha \leqslant \beta\) ）得出。

由这些评述可以推断，若诸广群 \(E_{\alpha}\) 是幺半群（相应地，群），且诸 \(f_{\alpha\beta}\) 是幺半群同态，则广群 E 是幺半群（相应地，群）。这时我们说这是一个幺半群（相应地，群）的逆向系统。泛性质立即推广到这种情形。

留给读者定义环的逆向系统 \((\mathrm{E}_{\alpha}, f_{\alpha \beta})\) 并验证 \(\mathbf{E} = \varprojlim \mathbf{E}_{\alpha}\) 是积环 \(\prod_{\alpha \in I} \mathbf{E}_{\alpha}\) 的子环，称为诸环 \(\mathbf{E}_{\alpha}\) 的逆极限环；可以验证，泛性质在此情形下仍然成立。

设 \(\mathfrak{E} = (\mathrm{E}_{\alpha}, f_{\alpha\beta})\) 与 \(\mathfrak{E}' = (\mathrm{E}_{\alpha}', f_{\alpha\beta}')\) 是广群（相应地，幺半群、群、环）的两个逆向系统。从 \(\mathfrak{E}\) 到 \(\mathfrak{E}'\) 的一个同态是一个映射的逆向系统 \((u_{\alpha})_{\alpha \in I}\) （ \(u_{\alpha}: \mathrm{E}_{\alpha} \to \mathrm{E}_{\alpha}'\) ），使得每个 \(u_{\alpha}\) 都是同态。在这些条件下，映射 \(u = \varprojlim u_{\alpha}\) 从 \(\varprojlim \mathrm{E}_{\alpha}\) 到 \(\varprojlim \mathrm{E}_{\alpha}'\) 是一个同态（参见《集合论》III，§ 7，第 2 号）。

\subsection{作用的逆向极限}

假设给定相对于同一指标集 I 的两个集合逆向系统 \((\Omega_{\alpha}, \phi_{\alpha\beta})\) 与 \((\mathrm{E}_{\alpha}, f_{\alpha\beta})\) 。假设对所有 \(\alpha \in I\) 给定了 \(\Omega_{\alpha}\) 在 \(E_{\alpha}\) 上的一个作用，使得

\[
f _ {\alpha \beta} (\omega_ {\beta} x _ {\beta}) = \phi_ {\alpha \beta} (\omega_ {\beta}) \cdot f _ {\alpha \beta} (x _ {\beta})\tag{1}
\]

对 \(\alpha \leqslant \beta, x_{\beta} \in \mathrm{E}_{\beta}, \omega_{\beta} \in \Omega_{\beta}\) 。那么所考虑的这族作用称为一个作用逆向系统。设 \(\Omega = \varprojlim \Omega_{\alpha}\) 与 \(\mathrm{E} = \varprojlim \mathrm{E}_{\alpha}\) ；若 \(x = (x_{\alpha})_{\alpha \in I}\) 属于 \(\mathrm{E}\) 且 \(\omega = (\omega_{\alpha})_{\alpha \in I}\) 属于 \(\Omega\) ，则 \(\omega \cdot x = (\omega_{\alpha} \cdot x_{\alpha})_{\alpha \in I}\) 由(1)属于 E。这样就定义了 \(\Omega\) 在 E 上的一个作用，称为诸 \(\Omega_{\alpha}\) 在诸 \(\mathrm{E}_{\alpha}\) 上作用的逆极限。

上述内容特别适用于 \(\Omega_{\alpha}\) 是幺半群且每个 \(\Omega_{\alpha}\) 在 \(E_{\alpha}\) 上的作用是一个运算的情形。那么这些运算的逆极限是该幺半群在 E 上的一个运算。

留给读者定义带算子的群的逆系统的逆向极限，并验证该极限是带算子的群。

\subsection{广群的正向系统}

在本节及下一节中，我们假设 I 是右有向的。

定义2. 关于指标集I的广群正向系统是一个集合的正向系统 \((\mathrm{E}_{\alpha}, f_{\beta\alpha})\) ，其中每个 \(E_{\alpha}\) 具有一个广群结构，且每个 \(f_{\beta\alpha}\) 都是广群同态。

设 \((\mathrm{E}_{\alpha}, f_{\beta\alpha})\) 是一个广群正向系统。用 E 表示正向极限集 \(\lim_{\rightarrow} \mathrm{E}_{\alpha}\) ，用 \(f_{\alpha}\) 表示从 \(\mathrm{E}_{\alpha}\) 到 E 的典范映射。回顾

\begin{enumerate}
\def\labelenumi{(\arabic{enumi})}
\setcounter{enumi}{1}
\tightlist
\item
\end{enumerate}

\[
f _ {\beta} \circ f _ {\beta \alpha} = f _ {\alpha} \quad \text { for } \alpha \leqslant \beta ,\tag{3}
\]

\[
\mathrm{E} = \bigcup_ {\alpha \in \mathbf {I}} f _ {\alpha} (\mathrm{E} _ {\alpha}).
\]

由(2)，同样地

\[
f _ {\alpha} (\mathrm{E} _ {\alpha}) \subset f _ {\beta} (\mathrm{E} _ {\beta}) \quad \text { for } \alpha \leqslant \beta .\tag{4}
\]

若 \(x_{\alpha}, y_{\alpha} \in \mathrm{E}_{\alpha}\) 满足 \(f_{\alpha}(x_{\alpha}) = f_{\alpha}(y_{\alpha})\) ，则存在 \(\beta \geqslant \alpha\) 使得 \(f_{\beta \alpha}(x_{\alpha}) = f_{\beta \alpha}(y_{\alpha})\) 。

命题1. 在 E 上存在一种且仅一种广群结构，使得映射 \(f_{\alpha} \colon \mathrm{E}_{\alpha} \to \mathrm{E}\) 是同态。若诸广群 \(\mathrm{E}_{\alpha}\) 是结合的（相应地，交换的），则 E 亦然。若诸广群 \(\mathrm{E}_{\alpha}\) 与诸同态 \(f_{\beta \alpha}\) 是保单位的，则广群 E 与诸同态 \(f_{\alpha}\) 亦然。

诸广群 \(E_{\alpha}\) 将以乘法记法书写。

设 \(x, y\) 是 E 中的元素。存在属于 I 的 \(\alpha\) 以及 \(x_{\alpha}, y_{\alpha}\) 属于 \(E_{\alpha}\) ，使得 \(x = f_{\alpha}(x_{\alpha})\) 与 \(y = f_{\alpha}(y_{\alpha})\) 。若 E 上存在使 \(f_{\alpha}\) 为同态的广群结构，则 \(x.y = f_{\alpha}(x_{\alpha}y_{\alpha})\) ，由此得该广群结构的唯一性。

为证明存在性，我们必须证明：对 I 中的 \(\alpha, \beta\) ， \(x_{\alpha}, y_{\alpha}\) 属于 E \(_{\alpha}\) 且 \(x_{\beta}', y_{\beta}'\) 属于 E \(_{\beta}\) 时，关系

\[
f _ {\alpha} (x _ {\alpha}) = f _ {\beta} (x _ {\beta} ^ {\prime}), \quad f _ {\alpha} (y _ {\alpha}) = f _ {\beta} (y _ {\beta} ^ {\prime})\tag{5}
\]

蕴含 \(f_{\alpha}(x_{\alpha}y_{\alpha}) = f_{\beta}(x_{\beta}'y_{\beta}')\) 。对于 \(\gamma \geqslant \alpha\) 和 \(\gamma \geqslant \beta\) ，我们设 \(x_{\gamma} = f_{\gamma \alpha}(x_{\alpha})\) , \(y_{\gamma} = f_{\gamma \alpha}(y_{\alpha}), x_{\gamma}^{\prime} = f_{\gamma \beta}(x_{\beta}^{\prime}), y_{\gamma}^{\prime} = f_{\gamma \beta}(y_{\beta}^{\prime})\) 。根据正向极限的定义，存在 I 中的 \(\gamma\) 满足 \(\gamma \geqslant \alpha, \gamma \geqslant \beta, x_{\gamma} = x_{\gamma}^{\prime}, y_{\gamma} = y_{\gamma}^{\prime}\) 。则

\[
\begin{array}{r l} f _ {\alpha} (x _ {\alpha} y _ {\alpha}) & = f _ {\gamma} (f _ {\gamma \alpha} (x _ {\alpha} y _ {\alpha})) = f _ {\gamma} (x _ {\gamma} y _ {\gamma}) = f _ {\gamma} (x _ {\gamma} ^ {\prime} y _ {\gamma} ^ {\prime}) = f _ {\gamma} (f _ {\gamma \beta} (x _ {\beta} ^ {\prime} y _ {\beta} ^ {\prime})) \\ & = f _ {\beta} (x _ {\beta} ^ {\prime} y _ {\beta} ^ {\prime}). \end{array}
\]

假设诸广群 \(E_{\alpha}\) 是结合的。设 x、y、z 属于 E。存在 \(\alpha \in I\) 及 \(x_{\alpha}, y_{\alpha}, z_{\alpha}\) 属于 \(E_{\alpha}\) ，使得

\[
x = f _ {\alpha} (x _ {\alpha}), \quad y = f _ {\alpha} (y _ {\alpha}), \quad z = f _ {\alpha} (z _ {\alpha}).
\]

则 \(xy = f_{\alpha}(x_{\alpha}y_{\alpha})\) ，从而 \((xy)z = f_{\alpha}((x_{\alpha}y_{\alpha})z_{\alpha})\) ；类似地

\[
x (y z) = f _ {\alpha} \left(x _ {\alpha} \left(y _ {\alpha} z _ {\alpha}\right)\right),
\]

从而 \((xy)z = x(yz)\) ，因为 \((x_{\alpha}y_{\alpha})z_{\alpha} = x_{\alpha}(y_{\alpha}z_{\alpha})\) 。交换广群的情形可类似地处理。

最后假设每个广群 \(E_{\alpha}\) 有单位元 \(e_{\alpha}\) ，且 \(f_{\beta\alpha}(e_{\alpha}) = e_{\beta}\) 对 \(\alpha \leqslant \beta\) 成立。对 I 中的 \(\alpha, \beta\) ，存在 I 中的 \(\gamma\) 满足 \(\gamma \geqslant \alpha\) 与 \(\gamma \geqslant \beta\) ，从而

\[
f _ {\alpha} (e _ {\alpha}) = f _ {\gamma} \left(f _ {\gamma \alpha} (e _ {\alpha})\right) = f _ {\gamma} (e _ {\gamma}) = f _ {\gamma} \left(f _ {\gamma \beta} (e _ {\beta})\right) = f _ {\beta} (e _ {\beta})
\]

因此存在元素 \(e\) 于 \(\mathbf{E}\) 中，使得 \(f_{\alpha}(e_{\alpha}) = e\) 对一切 \(\alpha \in \mathbf{I}\) 成立。设 \(x\) 是 \(\mathbf{E}\) 中的元素；选取 \(\alpha \in \mathbf{I}\) 与 \(x_{\alpha} \in \mathrm{E}_{\alpha}\) ，使得 \(x = f_{\alpha}(x_{\alpha})\) 。则

\[
e x = f _ {\alpha} (e _ {\alpha}) \cdot f _ {\alpha} (x _ {\alpha}) = f _ {\alpha} (e _ {\alpha}. x _ {\alpha}) = f _ {\alpha} (x _ {\alpha}) = x
\]

并且类似地有 \(x.e = x\) ，因此 \(e\) 是 \(E\) 的单位元。

广群 E 称为诸广群 \(E_{\alpha}\) 的正向极限。

命题2. 设 \((\mathrm{E}_{\alpha}, f_{\beta \alpha})\) 是一个广群正向系统，E 为其正向极限， \(f_{\alpha} \colon \mathrm{E}_{\alpha} \to \mathrm{E}\) 为典范同态。假设给定一个广群 F 及同态 \(u_{\alpha} \colon \mathrm{E}_{\alpha} \to \mathrm{F}\) ，使得 \(u_{\alpha} = u_{\beta} \circ f_{\beta \alpha}\) 对 \(\alpha \leqslant \beta\) 成立。存在一个且仅一个同态 \(u \colon \mathrm{E} \to \mathrm{F}\) ，使得 \(u_{\alpha} = u \circ f_{\alpha}\) 对所有 \(\alpha \in \mathbf{I}\) 成立。若诸广群 \(\mathrm{E}_{\alpha}\) 与 F 以及诸同态 \(f_{\beta \alpha}\) 与 \(u_{\alpha}\) 都是保单位的，则同态 \(u\) 是保单位的。

我们知道（《集合论》III，§ 7，第 6 号，命题6）存在一个且仅一个映射 \(u: \mathrm{E} \to \mathrm{F}\) ，使得 \(u_{\alpha} = u \circ f_{\alpha}\) 对所有 \(\alpha \in \mathbf{I}\) 成立。我们验证 \(u\) 是一个同态：设 \(x, y\) 属于 \(\mathrm{E}, \alpha\) 属于 \(\mathrm{I}\) ，且 \(x_{\alpha}, y_{\alpha}\) 属于 \(\mathrm{E}_{\alpha}\) ，满足 \(x = f_{\alpha}(x_{\alpha})\) 与 \(y = f_{\alpha}(y_{\alpha})\) 。则 \(xy = f_{\alpha}(x_{\alpha}y_{\alpha})\) ，从而

\[
\begin{array}{r l} u (x y) & = u (f _ {\alpha} (x _ {\alpha} y _ {\alpha})) = u _ {\alpha} (x _ {\alpha} y _ {\alpha}) = u _ {\alpha} (x _ {\alpha}) u _ {\alpha} (y _ {\alpha}) \\ & = u (f _ {\alpha} (u _ {\alpha})) u (f _ {\alpha} (y _ {\alpha})) = u (x) u (y). \end{array}
\]

现在考虑保单位的情形，并用 \(e_{\alpha}\) 表示 \(E_{\alpha}\) 的单位元， \(e\) 表示 \(E\) 的单位元， \(e'\) 表示 \(F\) 的单位元。选取 \(\alpha \in I\) ，则 \(e = f_{\alpha}(e_{\alpha})\) ，从而

\[
u (e) = u \left(f _ {\alpha} \left(e _ {\alpha}\right)\right) = u _ {\alpha} \left(e _ {\alpha}\right) = e ^ {\prime}
\]

因为 \(u_{\alpha}\) 是保单位的。因此 \(u\) 是保单位的。

与广群正向系统的概念类似，可以表述幺半群或群的正向系统的概念。命题1表明，作为幺半群正向系统 \((\mathrm{E}_{\alpha}, f_{\beta \alpha})_{\alpha, \beta \in I}\) 之极限的广群 E 是一个幺半群。我们证明，若诸 \(\mathrm{E}_{\alpha}\) 是群，则 E 是群：设 \(x \in \mathrm{E}\) , \(\alpha \in I\) 与 \(x_{\alpha} \in \mathrm{E}_{\alpha}\) 满足 \(x = f_{\alpha}(x_{\alpha})\) ；E 的元素 \(y = f_{\alpha}(x_{\alpha}^{-1})\) 是 \(x\) 的逆元（§ 2，第 3 号）。命题2的泛性质立即推广到幺半群或群的正向系统的情形。

留给读者定义环的正向系统。设 \((\mathbf{A}_{\alpha}, f_{\beta \alpha})\) 是这样一个正向系统；令 \(\mathbf{A} = \lim_{\longrightarrow} \mathbf{A}_{\alpha}\) ，并令 \(f_{\alpha}: \mathbf{A}_{\alpha} \to \mathbf{A}\) 为典范同态。在 \(\mathbf{A}\) 上存在（命题2）一种且仅一种加法与乘法，由 \(x + y = f_{\alpha}(x_{\alpha} + y_{\alpha})\) , \(xy = f_{\alpha}(x_{\alpha}y_{\alpha})\) （对 \(\alpha \in \mathbf{I}\) 、 \(x_{\alpha}, y_{\alpha}\) 属于 \(\mathbf{A}_{\alpha}\) 且 \(x = f_{\alpha}(x_{\alpha})\) , \(y = f_{\alpha}(y_{\alpha})\) ）刻画。在加法下 \(\mathbf{A}\) 是交换群，乘法是结合的且保单位的。最后，对 \(x, y, z\) 属于 \(\mathbf{A}\) ，选取 \(\alpha\) 属于 \(\mathbf{I}\) 及 \(x_{\alpha}, y_{\alpha}\) 与 \(z_{\alpha}\) 属于 \(\mathbf{A}_{\alpha}\) ，使得

\[
x = f _ {\alpha} (x _ {\alpha}), \quad y = f _ {\alpha} (y _ {\alpha}) \quad \text { and } \quad z = f _ {\alpha} (z _ {\alpha}).
\]

那么

\[
\begin{array}{r l} (x + y) \cdot z & = f _ {\alpha} (x _ {\alpha} + y _ {\alpha}) f _ {\alpha} (z _ {\alpha}) = f _ {\alpha} ((x _ {\alpha} + y _ {\alpha}) z _ {\alpha}) \\ & = f _ {\alpha} (x _ {\alpha} z _ {\alpha} + y _ {\alpha} z _ {\alpha}) = f _ {\alpha} (x _ {\alpha} z _ {\alpha}) + f _ {\alpha} (y _ {\alpha} z _ {\alpha}) = x z + y z \end{array}
\]

而关系 \(x(y + z) = xy + xz\) 可类似地证明。换言之，A 具有一个环结构，其特征由 \(f_{\alpha}\) 对所有 \(\alpha \in I\) 是环同态这一事实刻画。

环 A 称为诸环 \(A_{\alpha}\) 的正向极限。命题2立即推广到环的情形。

命题3. (a) 若诸 \(\mathbf{A}_{\alpha}\) 非零，则 \(\mathbf{A}\) 非零。

\begin{enumerate}
\def\labelenumi{(\alph{enumi})}
\setcounter{enumi}{1}
\item
  若诸 \(\mathbf{A}_{\alpha}\) 都是整环，则 \(\mathbf{A}\) 是整环。
\item
  若诸 \(\mathbf{A}_{\alpha}\) 都是域，则 \(\mathbf{A}\) 是域。
\end{enumerate}

设 \(0_{\alpha}, 1_{\alpha}\) 为 \(A_{\alpha}\) 的零元与单位元，0、1 为 A 的零元与单位元。存在 \(\alpha \in I\) 使得 \(f_{\alpha}(0_{\alpha}) = 0, f_{\alpha}(1_{\alpha}) = 1\) 。若 0 = 1，则存在 \(\beta \geqslant \alpha\) 使得 \(f_{\beta \alpha}(0_{\alpha}) = f_{\beta \alpha}(1_{\alpha})\) ，即 \(0_{\beta} = 1_{\beta}\) 。这就证明了 (a)。

假设诸 \(A_{\alpha}\) 是整环。则由 (a)，A 是交换的且非零。设 x、y 为 A 中满足 xy = 0 的元素。存在 \(\alpha \in I\) 与 \(x_{\alpha}, y_{\alpha} \in A_{\alpha}\) ，使得 \(x = f_{\alpha}(x_{\alpha}), y = f_{\alpha}(y_{\alpha})\) 。则 \(f_{\alpha}(x_{\alpha}y_{\alpha}) = xy = 0 = f_{\alpha}(0_{\alpha})\) 。因此存在 \(\beta \geqslant \alpha\) 使得 \(f_{\beta\alpha}(x_{\alpha}y_{\alpha}) = f_{\beta\alpha}(0_{\alpha})\) 。由于 \(A_{\beta}\) 是整环，由此得 \(f_{\beta\alpha}(x_{\alpha}) = 0_{\beta}\) 或 \(f_{\beta\alpha}(y_{\alpha}) = 0_{\beta}\) ，从而 x = 0 或 y = 0。这就证明了 (b)。

假设诸 \(A_{\alpha}\) 是域。则由 (a)， \(A \neq \{0\}\) 。设 x 是 A 的一个非零元素。存在 \(\alpha \in I\) 与 \(x_{\alpha} \in A_{\alpha}\) ，使得 \(x = f_{\alpha}(x_{\alpha})\) 。则 \(x_{\alpha} \neq 0\) 且 \(f_{\alpha}(x_{\alpha}^{-1})\) 是 x 在 A 中的逆元。这就证明了 (c)。

设 \(\mathfrak{E} = (\mathrm{E}_{\alpha}, f_{\beta \alpha})\) 与 \(\mathfrak{E}' = (\mathrm{E}_{\alpha}', f_{\beta \alpha}')\) 是广群（相应地，幺半群、群、环）的两个正向系统。从 \(\mathfrak{E}\) 到 \(\mathfrak{E}'\) 的一个同态是一个映射的正向系统 \((u_{\alpha})_{\alpha \in I}\) （ \(u_{\alpha} \colon E_{\alpha} \to E_{\alpha}'\) ），使得每个 \(u_{\alpha}\) 都是同态。在这些条件下，映射 \(u = \lim u_{\alpha}\) 从 \(E = \lim E_{\alpha}\) 到 \(E' = \lim E_{\alpha}'\) 是一个同态（参见《集合论》III，§ 7，第 6 号）。

\subsection{作用的正向极限}

假设给定相对于同一指标集 I 的两个集合正向系统 \((\Omega_{\alpha}, \phi_{\beta\alpha})\) 与 \((\mathrm{E}_{\alpha}, f_{\beta\alpha})\) ，并且对每个 \(\alpha \in I\) 给定 \(\Omega_{\alpha}\) 在 \(E_{\alpha}\) 上的一个作用。假设

\[
f _ {\beta \alpha} (\omega_ {\alpha}. x _ {\alpha}) = \phi_ {\beta \alpha} (\omega_ {\alpha}). f _ {\beta \alpha} (x _ {\alpha})
\]

对 \(\alpha \leqslant \beta\) 、 \(\omega_{\alpha} \in \Omega_{\alpha}\) 与 \(x_{\alpha} \in \mathrm{E}_{\alpha}\) 。那么所考虑的这族作用称为一个作用正向系统。如同命题2那样可以验证，存在作用 \(h\) —— \(\Omega = \varinjlim \Omega_{\alpha}\) 在 \(\mathrm{E} = \varinjlim \mathrm{E}_{\alpha}\) 上的作用——其描述如下：设 \(\omega \in \Omega\) 与 \(x \in \mathrm{E}\) ；选取 \(\alpha \in \mathrm{I}\) 及 \(\omega_{\alpha} \in \Omega_{\alpha}\) , \(x_{\alpha} \in \mathrm{E}_{\alpha}\) ，使得 \(\omega = \phi_{\alpha}(\omega_{\alpha})\) 与 \(x = f_{\alpha}(x_{\alpha})\) （ \(\phi_{\alpha}: \Omega_{\alpha} \to \Omega\) 与 \(f_{\alpha}: \mathrm{E}_{\alpha} \to \mathrm{E}\) 表示典范映射）；则 \(\omega . x = f_{\alpha}(\omega_{\alpha}. x_{\alpha})\) 。 \(\Omega\) 在 \(\mathrm{E}\) 上的这个作用称为诸 \(\Omega_{\alpha}\) 在诸 \(\mathrm{E}_{\alpha}\) 上作用的正向极限。

若诸 \(\Omega_{\alpha}\) 是幺半群，且每个 \(\Omega_{\alpha}\) 在 \(E_{\alpha}\) 上的作用是一个运算，则正向极限作用是一个运算。

留给读者定义带算子的群的正向系统的正向极限，并验证该极限是带算子的群。

\BourbakiExercisesSection

\begin{enumerate}
\def\labelenumi{\arabic{enumi}.}
\tightlist
\item
  设 E 是一个集合， \(A \subset E \times E\) 与 \((x, y) \mapsto x \top y\) , \((x, y) \in A\) 是 E 上的一个未必处处有定义的内部合成律。给定 E 的两个子集 X 与 Y，令 \(X \top Y\) 表示 E 的形如 \(x \top y\) （其中 \(x \in X\) , \(y \in Y\) 且 \((x, y) \in A\) ）的元素所成的集合。这样就在 \(\mathfrak{P}(E)\) 上定义了一个处处有定义的合成律。
\end{enumerate}

若 \((\mathbf{X}_{\alpha})_{\alpha \in \mathbf{A}}\) 与 \((\mathbf{Y}_{\beta})_{\beta \in B}\) 是 E 的子集的任意两个族，则

\[
\left(\bigcup_ {\alpha \in A} X _ {\alpha}\right) \top \left(\bigcup_ {\beta \in B} Y _ {\beta}\right) = \bigcup_ {(\alpha , \beta) \in A \times B} (X _ {\alpha} \top Y _ {\beta}).
\]

\begin{enumerate}
\def\labelenumi{\arabic{enumi}.}
\setcounter{enumi}{1}
\item
  设 \(\top\) 是集合 E 上未必处处有定义的一个合成律。设 E' 为 \(\mathfrak{P}(E)\) 的由诸集合 \(\{x\}\) （其中 \(x\) 遍历 E）与 E 的空子集 \(\varnothing\) 所组成的子集；证明 E' 是 \(\mathfrak{P}(E)\) 在律 \((X,Y)\mapsto X\top Y\) （习题1）下的稳定子集；由此推出，若 \(\overline{E}\) 表示通过向 E 添加（《集合论》II，§ 4，第 8 号）一个元素 \(\omega\) 而得到的集合，则律 \(\top\) 可以延拓到 \(\overline{E} \times \overline{E}\) 上，使得律 \(\top\) 与该延拓律在 E 上诱导的律相同。
\item
  设 \(\top\) 是 E 上并非处处有定义的一个合成律。
\end{enumerate}

\begin{enumerate}
\def\labelenumi{(\alph{enumi})}
\item
  对于律 \((\mathbf{X}, \mathbf{Y}) \mapsto \mathbf{X} \top \mathbf{Y}\) （习题1）—— \(\mathbf{E}\) 的子集之间的律——要使它是结合的，必要且充分的是：对所有 \(x \in \mathbb{E}\) , \(y \in \mathbb{E}\) , \(z \in \mathbb{E}\) ，若公式 \((x \top y) \top z = x \top (y \top z)\) 的两侧中有一侧有定义，则另一侧也有定义且与之相等（利用习题1证明该条件是充分的）。
\item
  若此条件成立，证明第 3 节定理 1 可推广如下：若公式 (3) 的两侧中有一侧有定义，则另一侧有定义且与之相等。
\end{enumerate}

\begin{enumerate}
\def\labelenumi{\arabic{enumi}.}
\setcounter{enumi}{3}
\tightlist
\item
  \begin{enumerate}
  \def\labelenumii{(\alph{enumii})}
  \tightlist
  \item
    给定集合 E，设 \(\Phi\) 为 E 的任意子集到 E 中的映射所成的集合；若 f、g、h 是 \(\Phi\) 的三个元素，证明若复合 \((f \circ g) \circ h\) 有定义，则 \(f \circ (g \circ h)\) 也有定义，但反之不成立；若这两个复合都有定义，则它们相等。
  \end{enumerate}
\end{enumerate}

\begin{enumerate}
\def\labelenumi{(\alph{enumi})}
\setcounter{enumi}{1}
\tightlist
\item
  设 \(\mathfrak{F}\) 是 \(E\) 的一族非空子集，其中任意两个都无公共元素，并令 \(\Psi\) 为 \(\Phi\) 的由 \(\mathfrak{F}\) 中一个集合到 \(\mathfrak{F}\) 中一个集合上的双射映射所组成的子集。证明：在由律 \(f \circ g\) 在 \(\Psi\) 上诱导的律下，习题3(a)的条件成立。
\end{enumerate}

\begin{enumerate}
\def\labelenumi{\arabic{enumi}.}
\setcounter{enumi}{4}
\item
  证明自然数的三元组 \((m, n, p)\) （ \(\neq 0\) ）中，使得 \((m^n)^p = m^{n^p}\) 的只有： \((1, n, p)\) ，其中 \(n\) 与 \(p\) 任意； \((m, n, 1)\) ，其中 \(m\) 与 \(n\) 任意；以及 \((m, 2, 2)\) ，其中 \(m\) 任意。
\item
  设 \(\top\) 是集合 E 上的一个合成律；设 A 为 E 的由这样的元素 \(x\) 组成的子集：满足 \(x \top (y \top z) = (x \top y) \top z\) （对一切 \(y \in E\) , \(z \in E\) ）；证明 A 是 E 的稳定子集，且 \(\top\) 在 A 上诱导的律是结合的。
\item
  若 \(\top\) 是 \(\mathbf{E}\) 上的结合律， \(a\) 与 \(b\) 是 \(\mathbf{E}\) 的两个元素，则集合 \(\{a\} \top \mathbf{E}, \mathbf{E} \top \{b\}, \{a\} \top \mathbf{E} \top \{b\}, \mathbf{E} \top \{a\} \top \mathbf{E}\) 是 \(\mathbf{E}\) 在律 \(\top\) 下的稳定子集。
\item
  设 \(\top\) 是 E 上的一个结合律， \(a\) 是 E 的一个元素；对一切 \(x \in \mathrm{E}\) , \(y \in \mathrm{E}\) ，我们记 \(x \perp y = x \top a \top y\) ：证明律 \(\bot\) 是结合的。
\item
  在集合 E 上，映射 \((x,y)\mapsto x\) 与 \((x,y)\mapsto y\) 是互为相反的结合律。
\item
  设 X 与 Y 是集合 E 的任意两个子集，令 \(X \top Y = X \cup Y\) 若 \(X \cap Y = \varnothing\) ，令 \(X \top Y = E\) 若 \(X \cap Y \neq \varnothing\) ；证明这样在 \(\mathfrak{P}(E)\) 上定义的合成律是结合的且交换的。
\item
  设 \(\top\) 是 \(\mathbf{E}\) 上的一个结合律， \(\mathbf{A}\) 与 \(\mathbf{B}\) 是 \(\mathbf{E}\) 的在该律下稳定的两个子集；证明，若 \(\mathbf{B} \top \mathbf{A} \subset \mathbf{A} \top \mathbf{B}\) ，则 \(\mathbf{A} \top \mathbf{B}\) 是 \(\mathbf{E}\) 的稳定子集。
\item
  仅有的不同的自然数 \(\neq 0\) ，它们在律 \((x,y)\mapsto x^{y}\) 下可交换的，是 2 和 4。
\item
  证明在律 \((\mathbf{X}, \mathbf{Y}) \mapsto \mathbf{X} \circ \mathbf{Y}\) （ \(E \times E\) 的子集之间）下，中心是由 \(\varnothing\) 与对角线组成的集合。
\item
  证明在 E 到 E 的映射之间的合成律 \(f \circ g\) 下，中心由恒等映射组成。
\item
  律 \(\top\) （在集合 \(E\) 上）称为幂等的，如果 \(E\) 的所有元素在该律下都是幂等的（第 4 号），即若 \(x \top x = x\) 对一切 \(x \in E\) 成立。证明，若律 \(\top\) 在 \(E\) 上是结合的、交换的且幂等的，则关系 \(x \top y = y\) 是 \(E\) 上的一个序关系；若记 \(x \leqslant y\) ，则任意两个元素 \(x, y\) （属于 \(E\) ）都有一个最小上界（在该序关系下）等于 \(x \top y\) 。试求其逆。
\item
  设 \(\top\) 是集合 E 上的一个结合且交换的合成律。设 \(n\) 为整数 \(>0\) 。映射 \(x \mapsto \frac{n}{\top} x\) 是 E 到 E 中的一个态射。
\end{enumerate}

\BourbakiExercisesSection

\begin{enumerate}
\def\labelenumi{\arabic{enumi}.}
\item
  设 \(\top\) 是集合 E 上的一个合成律。用 F 表示 E 与由单个元素组成的集合 \(\{e\}\) 的和集（《集合论》II，§ 4，第 8 号），并把 E 与 \(\{e\}\) 和 F 中相应的子集等同起来，证明可以以一种且仅一种方式在 F 上定义一个合成律 \(\overline{\top}\) ，它在 E 上诱导出律 \(\top\) ，且在该律下 \(e\) 是单位元；若 \(\top\) 是结合的，则律 \(\overline{\top}\) 是结合的。（称 F 是由 E “通过添加一个单位元”而导出的。）
\item
  设 \(\mathsf{T}\) 是 E 上一个处处有定义的律。
\end{enumerate}

\begin{enumerate}
\def\labelenumi{(\alph{enumi})}
\item
  要使 \(\top\) 是结合的，必要且充分的是：每个左平移 \(\gamma_{x}\) 与每个右平移 \(\delta_{y}\) 在 E 到 E 的映射所成的集合中（配备律 \(f \circ g\) ）都可交换。
\item
  假设 E 有单位元。要使 T 是结合的且交换的，必要且充分的是：映射 \((x,y)\mapsto x\top y\) 是广群 \(E\times E\) 到广群 E 中的一个态射。
\end{enumerate}

\begin{enumerate}
\def\labelenumi{\arabic{enumi}.}
\setcounter{enumi}{2}
\item
  在 E 到 E 的映射所成的集合 F 中，关系 \(f \circ g = f \circ h\) 蕴含 g = h，必要且充分的是 f 是单射；关系 \(g \circ f = h \circ f\) 蕴含 g = h，必要且充分的是 f 是满射；f （在律 \(\circ\) 下）是可消去的，必要且充分的是 f 是双射。
\item
  要使集合E上存在一个合成法则，使得E的每一个置换都是在该法则下E到自身的同构，其充分必要条件是E含有0、1或3个元素。
\item
  对 \(2 \leqslant n \leqslant 5\) ，在一个具有 n 个元素的集合 E 上，确定所有处处有定义、具有单位元且使 E 的所有元素都可消去的律；对 n = 5，证明存在满足这些条件的非结合律。
\end{enumerate}

（注意：习题 6 至 17 涉及以乘法记法书写的结合律； \(e\) 表示单位元（当它存在时）， \(\gamma_a\) 与 \(\delta_a\) 表示由元素 \(a \in E\) 所作的平移；对每个子集 \(X\) （属于 \(E\) ），我们记 \(\gamma_a(X) = aX\) , \(\delta_a(X) = Xa\) 。）

\begin{enumerate}
\def\labelenumi{\arabic{enumi}.}
\setcounter{enumi}{5}
\item
  对于有限集合上的一个结合律，每个可消去元素都是可逆的（利用第 3 号，命题 3）。
\item
  给定集合 E 上的一个结合律与一个元素 \(x \in E\) ，设 A 为由 \(x^{n}\) （ \(n \in N^{*}\) ）组成的集合；若存在单位元，设 B 为由 \(x^{n}\) （ \(n \in N\) ）组成的集合；若进一步 x 可逆，设 C 为由 \(x^{a}\) （ \(a \in Z\) ）组成的集合。证明，若 A（相应地
\end{enumerate}

B，C）是无限的，则它（配备由 E 上给定的律所诱导的律）同构于带加法的 \(N^{*}\) （相应地，N，Z）。

\begin{enumerate}
\def\labelenumi{\arabic{enumi}.}
\setcounter{enumi}{7}
\item
  在习题7的记号中，设 A 是有限的；证明 A 包含一个且仅一个幂等元（§ 1，第 4 号）（注意，若 \(x^p\) 与 \(x^q\) 是幂等元，则 \(x^p = x^{pq}\) , \(x^q = x^{pq}\) ，因此 \(x^p = x^q\) ）；若 \(h = x^p\) ，则由诸 \(x^n\) （ \(n \geqslant p\) ）所成的集合是 E 的一个稳定子集 D，使得在 D 上诱导的律下，D 是一个群。
\item
  在集合 E 上的一个乘法律下，设 \(a\) 是 E 的一个左可消去的元素。
\end{enumerate}

\begin{enumerate}
\def\labelenumi{(\alph{enumi})}
\item
  若存在元素 \(u\) 使得 \(au = a\) ，证明 \(ux = x\) 对所有 \(x \in \mathbb{E}\) 成立；特别地，若 \(xu = x\) 对所有 \(x \in \mathbb{E}\) 成立，则 \(u\) 是单位元。
\item
  若存在 \(u \in E\) 使得 au = a，且存在 \(b \in E\) 使得 ab = u，证明 ba = u（考察 aba）；特别地，若存在单位元 e 及元素 b 使得 ab = e，则 b 是 a 的逆元。
\end{enumerate}

\begin{enumerate}
\def\labelenumi{\arabic{enumi}.}
\setcounter{enumi}{9}
\tightlist
\item
  若 a 与 b 是 E 的两个元素，且 ba 是左可消去的，证明 a 是左可消去的。由此推出，在 E 上的一个结合且交换的律下，E 的不可消去元素所成的集合 S 满足 ES ⊆ S（从而是稳定的）。
\end{enumerate}

¶ 11. 若 E 的每个元素 x 都是左可消去的，则称 E 为左半群。(a) 若 u 是幂等元（§ 1，第 4 号），则对所有 \(x \in E\) 有 ux = x（利用习题9(a)）；u 是 Eu 上诱导的律的单位元。

\begin{enumerate}
\def\labelenumi{(\alph{enumi})}
\setcounter{enumi}{1}
\item
  若 \(u\) 与 \(v\) 是 \(\mathbf{E}\) 的两个不同的幂等元，则 \(\mathrm{Eu} \cap \mathrm{Ev} = \varnothing\) ，且集合 \(\mathrm{Eu}\) 与 \(\mathrm{Ev}\) （配备由 \(\mathbf{E}\) 上的律诱导的律）是同构的。
\item
  设 R 为诸集合 Eu（u 遍历 E 的幂等元集合）的并集的补集。证明 ER ⊂ R；由此 R 是 E 的稳定子集。若 R 非空，则 \(aR \neq R\) 对所有 \(a \in R\) 成立（利用习题 9(a) 证明：否则 R 将包含一个幂等元）；特别地，此时 R 是无限的。R 称为左半群 E 的剩余。
\item
  若 R 非空且 E 中至少存在一个幂等元 u（即 \(E \neq R\) ），则对所有 \(x \in REu\) , \(xEu \neq Eu\) ；特别地，REu 中没有元素在 Eu 中可逆，且 REu 是无限集（利用习题 8）。
\item
  若 E 有单位元 e，则 e 是 E 中唯一的幂等元且 R 为空（注意 E = Ee）。
\item
  若 E 中存在一个右可消去的元素 \(a\) （特别地，若 E 上给定的律是交换的），则要么 E 有单位元，要么 \(\mathbf{E} = \mathbf{R}\) （注意，若存在幂等元 \(u\) ，则 \(xa = xua\) 对所有 \(x \in \mathbf{E}\) 成立）。例如，若 E 是整数 \(\geqslant 1\) 的集合并赋予加法，则 E 是一个交换半群且 \(\mathbf{E} = \mathbf{R}\) 。
\item
  若存在 \(a \in \mathbf{E}\) 使得 \(\delta_a\) 是满射的，则要么 \(\mathbf{E}\) 有单位元，要么 \(\mathbf{E} = \mathbf{R}\) （分别考察 \(a \in \mathbf{R}\) 的情形与 \(a \in \mathbf{Eu}\) （对某个幂等元 \(u\) ）的情形）。
\end{enumerate}

¶ 12. 在 E 上的一个乘法律下，设 \(a \in E\) 使得 \(\gamma_{a}\) 是满射的。

\begin{enumerate}
\def\labelenumi{(\alph{enumi})}
\item
  证明，若存在 \(u\) 使得 \(ua = a\) ，则 \(ux = x\) 对所有 \(x \in E\) 成立。
\item
  要使元素 \(b \in E\) 满足 ba 是左可消去的，必要且充分的是 \(\gamma_{a}\) 是满射的且 b 是左可消去的。
\end{enumerate}

¶ 13. 假设对所有 \(x \in E\) ， \(\gamma_x\) 都是满射的。证明若对某个元素 \(a \in E\) ， \(\gamma_a\) 是双射映射，则 \(\gamma_x\) 对所有 \(x \in E\) 都是双射映射（利用习题12(b)）；特别地，若存在两个元素 \(a, b\) （属于 \(E\) ）使得 \(ab = b\) （利用习题12(a)），则 \(E\) 是一个剩余 \(R\) 为空的左半群；此外，对每个幂等元 \(u\) ， \(Eu\) 的每个元素在 \(Eu\) 中可逆。

¶ 14. 对于有限集合 E 上的一个结合律，存在形如 aE 的极小子集（即按包含序排列的、具有这种形式的 E 的子集族中的极小元素）。

\begin{enumerate}
\def\labelenumi{(\alph{enumi})}
\item
  若 M = aE 是极小的，则对所有 \(x \in M\) 有 xM = xE = M；在由 E 上的律诱导的律下，M 是一个左半群（习题 11），其中每个左平移都是双射的（参见习题 13）。
\item
  若 \(\mathbf{M} = a\mathbf{E}\) 与 \(\mathbf{M}' = a'\mathbf{E}\) 是极小的且互不相同，则 \(\mathbf{M} \cap \mathbf{M}' = \varnothing\) ；对所有 \(b \in \mathbf{M}\) ，映射 \(x' \mapsto bx'\) （从 \(\mathbf{M}'\) 到 \(\mathbf{M}\) 中）是 \(\mathbf{M}'\) 到 \(\mathbf{M}\) 上的一个双射映射。由此推出，存在幂等元 \(u' \in \mathbf{M}'\) 使得 \(bu' = b\) ，以及幂等元 \(u \in \mathbf{M}\) 使得 \(uu' = u\) （取 \(u\) 为满足 \(bu = b\) 的幂等元）。证明 \(u'u = u'\) （考察 \(uu'u\) ），且每个 \(y' \in \mathbf{M}'\) 若满足 \(y'u = y'\) 便属于 \(\mathbf{M}'u'\) 。
\item
  证明映射 \(x' \mapsto ux'\) （从 \(M'u'\) 到 \(M\) 中）是 \(M'u'\) 到 \(Mu\) 上的一个同构；由此推出 \(M\) 与 \(M'\) 是同构的左半群。
\item
  设 \(M_{i}\) ( \(1 \leqslant i \leqslant r\) ) 是形如 aE 的各不相同的极小子集：由 (b) 推出，每个左半群 \(M_{i}\) 的幂等元可以排成一个序列 \((u_{ij})\) ( \(1 \leqslant j \leqslant s\) )，使得对所有 i、j、k 有 \(u_{ij}u_{kj} = u_{ij}\) 。若 K 是诸 \(M_{i}\) 的并集，证明 \(Eu_{ij} \subset K\) （注意，对所有 \(x \in E\) ， \(xu_{ij}E\) 是形如 aE 的极小子集）；由此推出 \(Eu_{ij}\) 是诸 \(u_{kj}Eu_{kj}\) （ \(1 \leqslant k \leqslant r\) ）的并集（证明 \((\mathrm{Eu}_{ij}) \cap \mathbf{M}_{k} = u_{kj}\mathrm{Eu}_{kj}\) ），且 \(Eu_{ij}E = K\) 。最后证明，形如 Eb 的每个极小子集都与这 s 个集合 \(Eu_{ij}\) 之一相同（注意，利用 (a)， \(Ebu_{ij}b = Eb\) ，并由此推出 \(Eb \subset K\) ）。
\end{enumerate}

\begin{enumerate}
\def\labelenumi{\arabic{enumi}.}
\setcounter{enumi}{14}
\tightlist
\item
  设 \((x_{i})_{1\leqslant i\leqslant n}\) 是左可消去元素的有限序列。
\end{enumerate}

\begin{enumerate}
\def\labelenumi{(\alph{enumi})}
\item
  证明关系 \(x_{1}x_{2}\ldots x_{n} = e\) 蕴含所有关系 \(x_{i + 1}\ldots x_nx_1x_2\ldots x_i = e\) （它们彼此由“循环置换”得到），对 \(1\leqslant i\leqslant n\) 。
\item
  由此推出，若序列 \((x_{i})\) 的复合是可逆的，则每个 \(x_{i}\) 都是可逆的。
\end{enumerate}

¶ 16. 设 \(\top\) 是集合 E 上的一个结合律，E* 为 E 的可消去元素所成的集合；假设 E* 非空且每个可消去元素都是中心元素。设 \(\mathfrak{F}\) 表示具有以下性质的子集 \(\mathbf{X}\) （ \(\mathbf{E}\) 的）所成的集合：存在 \(y \in \mathrm{E}^*\) 使得 \(\delta_y(\mathrm{E}) \subset \mathbf{X}\) 。

\begin{enumerate}
\def\labelenumi{(\alph{enumi})}
\item
  证明 \(\mathfrak{F}\) 中两个集合的交集属于 \(\mathfrak{F}\) 。
\item
  设 \(\Phi\) 为这样的函数 \(f\) 所成的集合：定义在 \(\mathfrak{F}\) 中一个集合上，取值于
\end{enumerate}

E 且满足：对 \(X \in \mathfrak{F}\) ， \(f(\mathbf{X})\) 属于 F。设 R 表示 \(\Phi\) 的元素 f 与 g 之间的关系：“存在集合 \(X \in F\) 使得 f 与 g 在 X 上的限制相同”。证明 R 是一个等价关系；设 \(\Psi = \Phi/R\) 为 \(\Phi\) 在该关系下的商集。

\begin{enumerate}
\def\labelenumi{(\alph{enumi})}
\setcounter{enumi}{2}
\item
  设 \(f\) 与 \(g\) 是 \(\Phi\) 的两个元素， \(A \in \mathfrak{F}\) 与 \(B \in \mathfrak{F}\) 为 \(f\) 与 \(g\) 分别定义于其上的集合；存在 \(X \subset B\) 属于 \(\mathfrak{F}\) ，使得 \(g(X) \subset A\) ；若 \(g_X\) 是 \(g\) 在 \(X\) 上的限制，证明映射 \(f \circ g_X\) 属于 \(\Phi\) ，且其类（模 R）仅依赖于 \(f\) 与 \(g\) 的类（而不依赖于 \(X\) ）；这个类称为 \(f\) 的类与 \(g\) 的类的复合；证明这样在 \(\Psi\) 上定义的合成律是结合的且有单位元。
\item
  对所有 \(a \in \mathbb{E}\) ，证明左平移 \(\gamma_{a}\) 属于 \(\Phi\) ；令 \(\phi_{a}\) 为其类（模 R）。证明映射 \(x \mapsto \phi_{x}\) 是 E 到 \(\Psi\) 的一个子幺半群上的同构，且若 \(x \in \mathbb{E}^{*}\) ，则 \(\phi_{x}\) 在 \(\Psi\) 中可逆（考察 \(\gamma_{x}\) 的逆映射，证明它属于 \(\Phi\) 且其类（模 R）是 \(\phi_{x}\) 的逆元）。由此推出第 4 号定理 1 的一个推广。
\end{enumerate}

\begin{enumerate}
\def\labelenumi{\arabic{enumi}.}
\setcounter{enumi}{16}
\tightlist
\item
  设E是一个幺半群，S是E的一个稳定子集，满足以下性质：
\end{enumerate}

\((\alpha)\) 对所有 \(s \in S\) 以及所有 \(a \in E\) ，存在 \(t \in S\) 与 \(b \in E\) 使得 \(sb = at\) 。

(β) 对所有 \(a, b \in \mathbb{E}\) 以及所有 \(s \in \mathbb{S}\) （满足 \(sa = sb\) ），存在 \(t \in \mathbb{S}\) 使得 \(at = bt\) 。

\begin{enumerate}
\def\labelenumi{(\alph{enumi})}
\tightlist
\item
  在 E × S 中用 \((a, s) \sim (b, t)\) 记关系：“存在 S 中的 \(s'\) 与 \(t'\) 使得 \(tt' = ss'\) 且 \(as' = bt'\) 。”
\end{enumerate}

证明 \(\sim\) 是一个等价关系。我们记 \(\overline{\mathbf{E}} = \mathbf{E} \times \mathbf{S}/\sim\) ，用 \(as^{-1}\) 表示 \((a, s) (\text{mod } \sim)\) 的等价类，并用 \(\varepsilon: \mathbf{E} \to \overline{\mathbf{E}}\) 表示映射 \(a \mapsto ae^{-1}\) 。

\begin{enumerate}
\def\labelenumi{(\alph{enumi})}
\setcounter{enumi}{1}
\item
  证明在 \(\overline{\mathbf{E}}\) 上存在一种且仅一种幺半群结构，使得 \(\varepsilon\) 是保单位同态，且对所有 \(s \in \mathbf{S}\) ， \(\varepsilon(s)\) 是可逆的。
\item
  证明 \((\overline{\mathbf{E}},\varepsilon)\) 具有定理 1（第 4 号）中所述的泛性质。
\item
  证明 \(\varepsilon\) 是单射当且仅当S中的元素都是可消去的。
\end{enumerate}

\BourbakiExercisesSection

\begin{enumerate}
\def\labelenumi{\arabic{enumi}.}
\tightlist
\item
  设 \(\alpha \mapsto f_{\alpha}\) 是 \(\Omega\) 在 E 上的一个作用。设 F 为 \(\Omega\) 在 \(\mathbf{E}^{\mathbb{E}}\) 中的像，G 为 \(\mathbf{E}^{\mathbb{E}}\) 在律 \((f,g) \mapsto f \circ g\) 下由 F 与 E 的恒等映射生成的稳定子集。
\end{enumerate}

\begin{enumerate}
\def\labelenumi{(\alph{enumi})}
\item
  证明 E 的每个在 \(\Omega\) 的作用下稳定的子集，在 \(E^{E}\) 在 E 上的典范作用限制到 G 上之下也是稳定的（第 1 号，例 3）。
\item
  设 X 是 E 的一个子集；证明由 X 生成的 E 的稳定子集是诸 \(f(x)\) （f 遍历 G，x 遍历 X）所成的集合。
\end{enumerate}

\begin{enumerate}
\def\labelenumi{\arabic{enumi}.}
\setcounter{enumi}{1}
\item
  设 \(\bot\) 是 E 上的一个关于结合的内律 \(\top\) 双重分配的内律；证明，若 \(x \perp x'\) 与 \(y \perp y'\) 在律 \(\top\) 下是可消去的，则 \(x \perp y'\) 与 \(y \perp x'\) 在律 \(\top\) 下可交换（以两种不同的方式计算复合 \((x \top y) \perp (x' \top y')\) ）。特别地，若律 \(\bot\) 有单位元，则律 \(\top\) 下两个可消去的元素在该律下可交换；若 E 的所有元素在律 \(\top\) 下都可消去，则该律是交换的。
\item
  设 \(\top\) 与 \(\bot\) 是集合 \(E\) 上的两个内律，使得 \(\bot\) 关于 \(\top\) 是右分配的。
\end{enumerate}

\begin{enumerate}
\def\labelenumi{(\alph{enumi})}
\item
  若律 \(\top\) 有单位元 \(e\) ，则 \(x \perp e\) 在律 \(\top\) 下是幂等的，对所有 \(x \in \mathbb{E}\) ；若进一步存在 \(y\) 使得 \(x \perp y\) 在律 \(\top\) 下可消去，则 \(x \perp e = e\) 。
\item
  若律 \(\perp\) 有单位元 u，且存在 \(z \in E\) 对律 \(\perp\) 与 \(\top\) 都可消去，则 u 在律 \(\top\) 下可消去。
\end{enumerate}

¶ 4. 设 \(\top\) 与 \(\bot\) 是 E 上的两个内律，各有一个单位元；若由每个律导出的 E 在自身上的左作用关于另一个律是分配的，则 E 的每个元素在这两个律下都是幂等的（若 \(e\) 是 \(\top\) 的单位元， \(u\) 是 \(\bot\) 的单位元，先证明 \(e \perp e = e\) ，注意 \(e = e \perp (u \top e)\) ）。

\begin{enumerate}
\def\labelenumi{\arabic{enumi}.}
\setcounter{enumi}{4}
\tightlist
\item
  在集合 E 上假设给定了三个内合成律：一个加法（不必结合也不必交换）、一个乘法（不必结合）以及一个记作 \(\top\) 的律。假设乘法有单位元 \(e\) ，律 \(\top\) 关于乘法是右分配的，且律 \(\top\) 关于加法是左分配的。证明：若存在 \(x, y, z\) 使得 \(x \top z, y \top z\) 与 \((x + y) \top z\) 在乘法下可消去，则 \(e + e = e\) （利用习题3(a)）。
\end{enumerate}

¶ 6. E 上的一个内合成律 \(\top\) 称为在 E 上确定一个拟群结构，如果对所有 \(x \in E\) ，左平移与右平移 \(\gamma_x\) 与 \(\delta_x\) 都是 E 到其自身的双射映射。称一个拟群是分配的，如果律 \(\top\) 关于其自身是双重分配的。

\begin{enumerate}
\def\labelenumi{(\alph{enumi})}
\item
  确定一个具有 \(n\) 个元素的集合上的所有分配拟群结构，对 \(2 \leqslant n \leqslant 6\) 。
\item
  *证明有理数集 \(\mathbf{Q}\) 配备律 \((x,y)\mapsto\frac{1}{2}(x+y)\) 是一个分配拟群。*
\item
  在分配拟群 E 中，每个元素都是幂等的。由此推出，若 E 含有多于一个元素，则律 \(\top\) 既不能有单位元，也不能是结合的。
\item
  分配拟群 E 的右平移和左平移都是 E 的自同构。
\item
  若E是有限分配拟群，则E的每个稳定子集上诱导的结构都是分配拟群结构。
\item
  设 E 是一个分配拟群。若 R 是一个与律 T 左（相应地，右）相容（第 3 号）的等价关系，则模 R 的类都是 E 的稳定子集。若 E 是有限的，则所有这些类都可以通过左（相应地，右）平移从一个得到另一个；在同一条件下，若 R 与律 T 相容，则 E/R 上的商结构是一个分配拟群结构。
\item
  设 E 是一个分配拟群。E 中与给定元素 a 可交换的元素所成的集合 \(A_{a}\) 是稳定的；若 E 是有限的，则对所有 \(x \in A_{a}\) 有 \(A_{x} = A_{a}\) （注意，利用 (e)，存在 \(y \in A_{a}\) 使得 \(x = a \top y\) ），且当 x 遍历 E 时，诸集合 \(A_{x}\) 与关于一个和律 T 相容的关系的等价类相同。
\item
  若 E 是有限分配拟群且律 \(\top\) 是交换的，则 E 的元素个数是奇数（考察 E 的元素的有序对 \((x,y)\) ，满足 \(x\top y=y\top x=a\) ，其中 a 是给定的）。
\end{enumerate}

\begin{enumerate}
\def\labelenumi{\arabic{enumi}.}
\setcounter{enumi}{6}
\tightlist
\item
  假设在集合 E 上给定一个（结合且交换的）加法，在该加法下 E 的所有元素都可逆，以及一个（结合的）乘法，它关于加法 \(*\) 是双重分配的（换言之，一个环结构） \(_{*}\) ；记 \([x,y]=xy-yx\) ；律 \((x,y)\mapsto[x,y]\) 关于加法是双重分配的。要使 x 与 y 在乘法下可交换，必要且充分的是 \([x,y]=0\) ；我们有恒等式
\end{enumerate}

\[
[ x, y ] = - [ y, x ]; \quad [ x, [ y, z ] ] + [ y, [ z, x ] ] + [ z, [ x, y ] ] = 0
\]

（第二个恒等式称为“雅可比恒等式”）。第二个恒等式也可以写成

\[
[ x, [ y, z ] ] - [ [ x, y ], z ] = [ [ z, x ], y ]
\]

它表达了律 \([x, y]\) 的“结合性偏差”。

\begin{enumerate}
\def\labelenumi{\arabic{enumi}.}
\setcounter{enumi}{7}
\tightlist
\item
  在与习题7相同的假设下，设 \(x \top y = xy + yx\) ；则律 \(\top\) 是交换的且关于加法是双重分配的，但一般不是结合的。
\end{enumerate}

\begin{enumerate}
\def\labelenumi{(\alph{enumi})}
\item
  对所有 \(x \in \mathbb{E}\) ，证明 \(\top^{m+n} x = \left( \begin{array}{c} m \\ \top \end{array} x \right) \top \left( \begin{array}{c} n \\ \top \end{array} x \right)\) 。
\item
  若记 \([x, y, z] = (x \top y) \top z - x \top (y \top z)\) （律 \(\top\) 的结合性偏差），证明恒等式
\end{enumerate}

\[
[ x, y, z ] + [ z, y, x ] = 0
\]

\[
[ x, y, z ] + [ y, z, x ] + [ z, x, y ] = 0
\]

\[
[ x \top y, u, z ] = [ u, [ (x \top y), z) ] ]
\]

（记号 \([x, y]\) 含义与习题7中相同）

\[
[ x \top y, u, z ] + [ y \top z, u, x ] + [ z \top x, u, y ] = 0.
\]

¶ 9. 假设在 E 上给定一个（结合且交换的）加法，使得 E 的所有元素在该加法下都可逆，并给定一个乘法，该乘法非结合，但交换且对加法双重分配。进一步假设 \(n \in \mathbf{Z}\) , \(n \neq 0\) 和 \(nx = 0\) 蕴含 \(x = 0\) 在 E 中成立。证明：若记 \([x, y, z] = (xy)z - x(yz)\) ，则恒等式

\[
[ x y, u, z ] + [ y z, u, z ] + [ z x, u, y ] = 0
\]

成立，则 \(x^{m + n} = x^m x^n\) 对所有 \(x\) 成立（对 \(p\) 用归纳法证明：恒等式 \([x^q, y, x^{p - q}] = 0\) 对 \(1 \leqslant q < p\) 成立）。

\begin{enumerate}
\def\labelenumi{\arabic{enumi}.}
\setcounter{enumi}{9}
\tightlist
\item
  设 E 是一个交换幺半群，其律以加法记法书写，其单位元记为 0。设 \(\top\) 是 E 上的一个关于加法分配的内律，且满足 \(0 \top x = x \top 0 = 0\) （对所有 \(x \in \mathbb{E}\) ）。设 S 是 E 的一个在加法下及在由 \(\top\) 导出的每个外律下都稳定的子集；令 \(\overline{\mathbb{E}}\) 表示 E 关于 S 的加法差幺半群， \(\varepsilon\) 为 E 到 \(\overline{\mathbb{E}}\) 中的典范同态。则在 \(\overline{\mathbb{E}}\) 上存在一种且仅一种律 \(\overline{\top}\) ，它关于加法是分配的，且使得 \(\varepsilon\) 对律 \(\top\) 与 \(\overline{\top}\) 是同态。若律 \(\top\) 是结合的（相应地，交换的），则律 \(\overline{\top}\) 亦然。
\end{enumerate}

\BourbakiExercisesSection

\begin{enumerate}
\def\labelenumi{\arabic{enumi}.}
\tightlist
\item
  确定一个具有 \(n\) 个元素的集合上的所有群结构，对 \(2 \leqslant n \leqslant 6\) （参见 § 2，习题5）。确定这些群的子群与商群，以及它们的若尔当–赫尔德序列。
\end{enumerate}

¶ 2. (a) 集合 E 上的一个结合律 \((x,y)\mapsto xy\) 是群律，如果存在 \(e\in\mathrm{E}\) 使得对所有 \(x\in\mathrm{E}\) 有 \(ex=x\) ，且如果对每个 \(x\in\mathrm{E}\) 存在 \(x'\) 使得 \(x^{\prime}x=e\) （证明 \(xx^{\prime}=e\) ，考察复合 \(x^{\prime}xx^{\prime}\) ；由此推出 \(e\) 是单位元）。

\begin{enumerate}
\def\labelenumi{(\alph{enumi})}
\setcounter{enumi}{1}
\tightlist
\item
  证明如果对所有 \(x \in E\) ，左平移 \(\gamma_{x}\) 是 E 到 E 上的一个映射，并且如果存在某个 \(a \in E\) 使得右平移 \(\delta_{a}\) 是 E 到 E 上的一个映射，则同样的结论成立。
\end{enumerate}

\begin{enumerate}
\def\labelenumi{\arabic{enumi}.}
\setcounter{enumi}{2}
\item
  在群 G 中，每个有限非空稳定子集 H 都是 G 的子群（参见 § 2，习题 6）。
\item
  设 A 和 B 是群 G 的两个子群。
\end{enumerate}

\begin{enumerate}
\def\labelenumi{(\alph{enumi})}
\item
  证明包含 A 和 B 的最小子群（即由 \(\mathbf{A} \cup \mathbf{B}\) 生成的子群）与奇数个（任取的）元素的序列 \((x_{i})_{1 \leqslant i \leqslant 2n + 1}\) 的复合所成的集合相同，其中 \(x_{i} \in \mathbf{A}\) 对 \(i\) 为奇数， \(x_{i} \in \mathbf{B}\) 对 \(i\) 为偶数。
\item
  为使AB成为G的子群（此时它是由A∪B生成的子群），其充分必要条件为A与B可交换，即AB = BA。
\item
  若 A 与 B 可交换，且 C 是包含 A 的子群，则 A 与 \(B \cap C\) 可交换，且 \(\mathbf{A}(\mathbf{B} \cap \mathbf{C}) = \mathbf{C} \cap (\mathbf{AB})\) 。
\end{enumerate}

\begin{enumerate}
\def\labelenumi{\arabic{enumi}.}
\setcounter{enumi}{4}
\item
  若群G的一个子群的指数为2，则它在G中正规。
\item
  设 \((\mathbf{G}_{\alpha})\) 是群 \(\mathbf{G}\) 的一族正规子群，使得 \(\bigcap_{\alpha} \mathbf{G}_{\alpha} = \{e\}\) ；证明 \(\mathbf{G}\) 同构于乘积群 \(\prod_{\alpha} (\mathbf{G} / \mathbf{G}_{\alpha})\) 的一个子群。
\item
  若 G 是两个子群 A 与 B 的直积，且 H 是 G 的一个子群，使得 \(A \subset H\) ，则 H 是 A 与 \(H \cap B\) 的直积。
\item
  设 A 与 A' 是两个群，且 G 是 A × A' 的一个子群。我们记：
\end{enumerate}

\[
\mathrm{N} = \mathrm{G} \cap (\mathrm{A} \times \{e \}), \quad \mathrm{H} = \operatorname{pr} _ {1} (\mathrm{G})
\]

\[
\mathrm{N} ^ {\prime} = \mathrm{G} \cap (\{e \} \times \mathrm{A} ^ {\prime}), \quad \mathrm{H} ^ {\prime} = \operatorname{pr} _ {2} (\mathrm{G}).
\]

\begin{enumerate}
\def\labelenumi{(\alph{enumi})}
\tightlist
\item
  证明 N 在 H 中正规，且 N' 在 H' 中正规；定义同构
\end{enumerate}

\[
\mathrm{H} / \mathrm{N} \rightarrow \mathrm{G} / (\mathrm{N} \times \mathrm{N} ^ {\prime}) \rightarrow \mathrm{H} ^ {\prime} / \mathrm{N} ^ {\prime}.
\]

\begin{enumerate}
\def\labelenumi{(\alph{enumi})}
\setcounter{enumi}{1}
\tightlist
\item
  假设 H = A， \(H' = A'\) ，群 A 与 \(A'\) 具有有限长度，且 A 的若尔当–赫尔德序列的任一商群都不同构于 \(A'\) 的若尔当–赫尔德序列的任一商群。证明 \(G = A \times A'\) 。
\end{enumerate}

\begin{enumerate}
\def\labelenumi{\arabic{enumi}.}
\setcounter{enumi}{8}
\item
  设 G 是一个不只由单位元组成的群。假设存在 G 的一个有限子集 S，使得 G 由 S（相应地，由元素 \(gsg^{-1}, g \in G, s \in S\) ）生成。设 N 为 G 的异于 G 的子群（相应地，正规子群）所成的集合。证明 N 在包含序下是归纳的。由此推出存在 G 的一个正规子群 H，使得 G/H 是单群。
\item
  设 H 是群 G 的一个正规子群，且包含在 G 的中心内。证明：若 G/H 是单生成群，则 G 是交换群。
\item
  如果群 G 中除单位元外的所有元素阶均为 2，则 G 是交换的；若 G 有限，则其阶 n 是 2 的幂（通过对 n 进行归纳论证）。
\item
  设 G 是一个群，使得对一个固定的整数 \(n > 1\) ， \((xy)^n = x^n y^n\) 对所有 \(x \in G\) 、 \(y \in G\) 成立。若 \(G^{(n)}\) 表示由 \(x^n\) （其中 \(x\) 遍历 G）组成的集合，而 \(G_{(n)}\) 表示由 \(x \in G\) （满足 \(x^n = e\) ）组成的集合，证明 \(G^{(n)}\) 与 \(G_{(n)}\) 是 G 的正规子群；若 G 有限， \(G^{(n)}\) 的阶等于 \(G_{(n)}\) 的指数。证明对所有 G 中的 \(x, y\) ，也有 \(x^{1-n}y^{1-n} = (xy)^{1-n}\) ，并由此推出 \(x^{n-1}y^n = y^n x^{n-1}\) ；由此得出结论：G 中形如 \(x^{n(n-1)}\) 的元素所成的集合生成 G 的一个交换子群。
\item
  设 G 是一个群。若 S 是一个单群，则称 S 出现在 G 中，如果存在 G 的两个子群 H、H'，其中 H 是 H' 的正规子群，使得 H'/H 同构于 S。令 in(G) 表示出现在 G 中的单群的同构类集合。
\end{enumerate}

\begin{enumerate}
\def\labelenumi{(\alph{enumi})}
\item
  证明 in(G) = ∅ ⇔ G = \{e\}.
\item
  若 \(\mathbf{H}\) 是 \(\mathbf{G}\) 的子群，证明 \(\operatorname{in}(\mathbf{H}) \subset \operatorname{in}(\mathbf{G})\) ；此外，若 \(\mathbf{H}\) 是正规的，则
\end{enumerate}

\[
\operatorname{in} (\mathrm{G}) = \operatorname{in} (\mathrm{H}) \cup \operatorname{in} (\mathrm{G} / \mathrm{H}).
\]

\begin{enumerate}
\def\labelenumi{(\alph{enumi})}
\setcounter{enumi}{2}
\item
  设 \(\mathbf{G}_1\) 和 \(\mathbf{G}_2\) 是两个群。证明以下两个性质的等价性：
\item
  \(\operatorname{in}(\mathbf{G}_1) \cap \operatorname{in}(\mathbf{G}_2) = \varnothing\) ;
\end{enumerate}

\begin{enumerate}
\def\labelenumi{(\roman{enumi})}
\setcounter{enumi}{1}
\tightlist
\item
  \(\mathbf{G}_1\times \mathbf{G}_2\) 的每个子群都具有形式 \(\mathbf{H}_1\times \mathbf{H}_2\) ，其中 \(\mathbf{H}_1\subset \mathbf{G}_1\) ， \(\mathbf{H}_2\subset \mathbf{G}_2\) 。
\end{enumerate}

（使用习题8。）

*当 \(\mathbf{G}_1\) 和 \(\mathbf{G}_2\) 是有限群时，证明（i）和（ii）等价于：

\begin{enumerate}
\def\labelenumi{(\roman{enumi})}
\setcounter{enumi}{2}
\tightlist
\item
  \(\operatorname{Card}(\mathbf{G}_1)\) 和 \(\operatorname{Card}(\mathbf{G}_2)\) 是互素的。*
\end{enumerate}

\begin{enumerate}
\def\labelenumi{\arabic{enumi}.}
\setcounter{enumi}{13}
\tightlist
\item
  设G为一个群，H为G的一个子群。G的任意子集T，若它与模H的每个左陪集恰好交于一点，则称T为模H的左陪集的一个代表系（参见集合论，II，§6，第2号）；此条件等价于以下条件：
\end{enumerate}

映射 \((x,y)\mapsto xy\) 是 \(\mathbf{T}\times \mathbf{H}\) 到 \(\mathbf{G}\) 上的一个双射。

\begin{enumerate}
\def\labelenumi{(\alph{enumi})}
\item
  设 \(\mathbf{T}\) 是这样的一个代表系；对所有 \(g \in \mathbf{G}\) 与所有 \(t \in \mathbf{T}\) ，设 \(x(g, t) \in \mathbf{T}\) 与 \(y(g, t) \in \mathbf{H}\) 满足 \(gt = x(g, t)y(g, t)\) 。设 \(\mathbf{S}\) 是 \(\mathbf{G}\) 的一个生成 \(\mathbf{G}\) 的子集。证明元素 \(y(g, t), g \in \mathbf{S}, t \in \mathbf{T}\) 生成 \(\mathbf{H}\) 。（设 \(\mathbf{H}'\) 为 \(\mathbf{H}\) 的由这些元素生成的子群，证明 \(\mathbf{T}.\mathbf{H}'\) 在诸平移 \(\boldsymbol{\gamma}_g, g \in \mathbf{S}\) 下稳定，从而在诸 \(\boldsymbol{\gamma}_g, g \in \mathbf{G}\) 下也稳定，并由此推出 \(\mathbf{T}.\mathbf{H}' = \mathbf{G}\) ，从而 \(\mathbf{H}' = \mathbf{H}\) 。）
\item
  假设 (G:H) 是有限的。证明：G 能被有限子集生成当且仅当 H 能被有限子集生成。
\end{enumerate}

¶ 15. 设F是带算子群G的一组稳定子群；若F的每个子集在包含序下都有极大（相应地，极小）元，则称F满足极大条件（相应地，极小条件）。

假设群G的所有稳定子群构成的集合满足极小条件。

\begin{enumerate}
\def\labelenumi{(\alph{enumi})}
\item
  证明不存在 G 的同构于 G 且异于 G 的稳定子群（用反证法论证：表明该假设将蕴含存在 G 的稳定子群的一个严格递减无限序列）。
\item
  G 的不只由 e 组成的正规稳定子群所成集合的极小元素称为极小正规稳定子群。设 M 是 G 的一族极小正规稳定子群，S 是 G 的包含 M 中所有子群的最小稳定子群；证明 S 是 G 的有限个极小正规稳定子群的直积（令 \((\mathbf{M}_{n})\) 为 G 的属于 M 且满足下述条件的一列极小正规稳定子群： \(M_{n+1}\) 不包含在由 \(M_{1}, M_{2}, \ldots, M_{n}\) 的并集生成的稳定子群中；令 \(S_{k}\) 为由诸 \(M_{n}\) （指标 \(n \geqslant k\) ）的并集生成的稳定子群；证明从某个秩开始 \(S_{k+1} = S_{k}\) ，因而 \((\mathbf{M}_{n})\) 是一个有限序列；最后使用第 9 号，命题 15）。
\item
  如果G是一个无算子的群，证明G的每个极小正规子群M是有限个彼此同构的单群的直积（设N是M的极小正规子群；证明M是G中包含所有子群 \(aNa^{-1}\) （其中a遍历G）的最小子群，然后将(b)应用于群M）。
\end{enumerate}

\begin{enumerate}
\def\labelenumi{\arabic{enumi}.}
\setcounter{enumi}{15}
\tightlist
\item
  如果带算子的群 G 的稳定子群集合满足极大条件和极小条件（习题 15），则 G 有一个若尔当–赫尔德序列（对 G 的一个子群 H，考虑 H 的异于 H 的正规稳定子群所成集合的一个极大元素）。
\end{enumerate}

¶ 17. 设G是一个带算子的群；G的一个合成列（G \(_{i}\) ）称为正规的，如果所有G \(_{i}\) 都是G的正规稳定子群；一个正规列Σ称为主列，如果它是严格递减的，并且不存在异于Σ、比Σ更细且严格递减的正规列。

\begin{enumerate}
\def\labelenumi{(\alph{enumi})}
\item
  若 \((\mathbf{G}_{i})\) 和 \((\mathbf{H}_{j})\) 是 G 的两个正规列，证明存在两个等价的正规列，分别比 \((\mathbf{G}_{i})\) 和 \((\mathbf{H}_{j})\) 更细（通过考虑 G 的一个合适的算子域来应用 Schreier 定理）。通过将子群 \(G_{ij}^{\prime} = G_{i} \cap (G_{i+1}H_{j})\) 和 \(H_{ji}^{\prime} = H_{j} \cap (H_{j+1}G_{i})\) “插入”到序列 \((\mathbf{G}_{i})\) 和 \((\mathbf{H}_{j})\) 中，给出这个命题的第二个证明。
\item
  如果G有主列，则G的任意两个主列等价；对于每个严格递减的正规列 \(\Sigma\) ，存在一个比 \(\Sigma\) 更细的主列。由此推出，G有主列的充分必要条件是G的正规稳定子群集合满足极大条件和极小条件。
\item
  如果G是一个无算子的群，有主列，并且G的子群集合满足极小条件，则每个商群
\end{enumerate}

\(G_{i}/G_{i+1}\) 是有限个彼此同构的单子群的直积（使用习题15）。

¶ 18. (a) 设 G 是一个带算子的群，由一族 G 的单正规稳定子群 \((\mathrm{H}_{i})_{i\in\mathbf{I}}\) 的并集生成。证明存在 I 的一个子集 J，使得 G 是族 \((\mathrm{H}_{i})_{i\in\mathbf{J}}\) 的限制和（将 Zorn 引理应用于 I 的具有下述性质的子集 K 所成的集合：由 \(H_{i}\) （ \(i\in K\) ）的并集生成的子群是该族的限制和）。

\begin{enumerate}
\def\labelenumi{(\alph{enumi})}
\setcounter{enumi}{1}
\tightlist
\item
  设 A 是 G 的一个正规稳定子群。证明存在 I 的一个子集 \(J_{A}\) ，使得 G 是 A 与诸 \(H_{i}\) （ \(i \in J_{A}\) ）的限制和。由此推出 A 同构于诸 \(H_{i}\) 的一个子族的限制和。
\end{enumerate}

¶ 19. (a) 设 G 是一个群，使得 G 的每个异于 G 的正规子群都包含在 G 的一个异于 G 的直因子中。证明 G 是一族单子群的限制和。（设 K 为 G 的由所有单子群的并集生成的子群。假设 K ≠ G；若 x \ensuremath{\in} G−K，考虑 G 的一个正规子群 M，使其在包含 K 而不包含 x 的那些正规子群中是极大的。若 G = M' × S，其中 M' 包含 M 且 S ≠ \{e\} ，证明 x ∉ M' ，从而 M' = M；另一方面，若 N 是 S 的一个正规子群，证明 M × N = G 并由此得出 S 是单群，从而得到矛盾。借助习题 18 得出结论。）得出逆命题（参见习题 18）。

\begin{enumerate}
\def\labelenumi{(\alph{enumi})}
\setcounter{enumi}{1}
\tightlist
\item
  为使在群 G 中存在一族 \((\mathrm{N}_{\alpha})\) 正规单子群，使得诸 \(G/N_{\alpha}\) 都是单群且 \(\bigcap_{\alpha} N_{\alpha} = \{e\}\) ，必要且充分的是：对 G 的每个正规子群 \(N \neq \{e\}\) ，存在一个正规子群 \(N' \neq G\) ，使得 \(NN' = G\) 。（要看出该条件是必要的，考虑一个异于 e 的 \(x \in N\) 以及一个不包含 x 的 \(N_{\alpha}\) 。要看出该条件是充分的，对 G 中所有 \(x \neq e\) ，考虑由 x 生成的正规子群 N 以及一个正规子群 \(N' \neq G\) ，使得 \(NN' = G\) 。证明：若 G 的正规子群 M 在包含 \(N'\) 而不包含 x 的 G 的正规子群中是极大的，则它在所有正规子群 \(\neq G\) 所成的集合中也是极大的，并由此得出结论。）
\end{enumerate}

\begin{enumerate}
\def\labelenumi{\arabic{enumi}.}
\setcounter{enumi}{19}
\tightlist
\item
  \begin{enumerate}
  \def\labelenumii{(\alph{enumii})}
  \tightlist
  \item
  设 G 是一个群，是两个子群 A 和 B 的直积。设 \(C_{A}\) 为 A 的中心，并设 N 是 G 的一个正规子群，使得 \(N \cap A = \{1\}\) 。证明 \(N \subset C_{A} \times B\) 。
  \end{enumerate}
\end{enumerate}

\begin{enumerate}
\def\labelenumi{(\alph{enumi})}
\setcounter{enumi}{1}
\tightlist
\item
  设 \((\mathbf{S}_i)_{i\in I}\) 是非交换单群的一个有限族。证明 \(\prod_{i\in I}S_i\) 的每个正规子群都等于诸部分积 \(\prod_{i\in J}S_i\) 之一，其中 \(J\) 是 \(I\) 的一个子集。
\end{enumerate}

\begin{enumerate}
\def\labelenumi{\arabic{enumi}.}
\setcounter{enumi}{20}
\tightlist
\item
  设 G 是一个有限长度的群。证明以下性质的等价性：
\end{enumerate}

\begin{enumerate}
\def\labelenumi{(\alph{enumi})}
\item
  G 是一些彼此同构的单群的乘积。
\item
  G 的任何异于 \(\{1\}\) 与 G 的子群在 G 的所有自同构下都不是稳定的。
\end{enumerate}

¶ 22. 设 E 是一个拟群（§ 3，习题 6），以乘法记号书写，具有一个幂等元 e，并且满足 \((xy)(zt) = (xz)(yt)\) 对所有 x, y, z, t 成立。设 \(u(x)\) 表示 E 中满足 \(u(x)e = x\) 的元素， \(v(x)\) 表示 E 中满足 \(ev(x) = x\) 的元素；证明合成律 \((x, y) \mapsto u(x)v(y)\) 是 E 上的一个交换群律，e 在其下是单位元，映射 \(x \mapsto xe\) 与 \(y \mapsto ey\) 是该群结构的可交换的自同态，并且 \(xy = u(xe)v(ey)\) （先建立恒等式 \(e(xy) = (ex)(ey)\) , \((xy)e = (xe)(ye)\) , \(e(xe) = (ex)e\) ；然后注意关系 x = y、ex = ey 与 xe = ye 是等价的）。得出逆命题。

¶ 23. 考虑集合 E 上的一个内律，它并非处处有定义，以乘法记号书写，并满足以下条件：

\begin{enumerate}
\def\labelenumi{(\arabic{enumi})}
\item
  如果合成 \((xy)z, x(yz)\) 之一有定义，则另一个也有定义，并且两者相等；
\item
  若 \(x, x', y\) 使得 \(xy\) 与 \(x'y\) （相应地 \(yx\) 与 \(yx'\) ）都有定义且相等，则 \(x = x'\) ；
\item
  对所有 \(x \in \mathbf{E}\) ，存在三个元素 \(e_x, e_x'\) 与 \(x^{-1}\) ，使得 \(e_xx = x\) , \(xe_x' = x\) , \(x^{-1}x = e_x'\) ； \(e_x\) 称为 \(x\) 的左单位元， \(e_x'\) 称为 \(x\) 的右单位元， \(x^{-1}\) 称为 \(x\) 的逆（这是语言上的一种滥用）。
\end{enumerate}

\begin{enumerate}
\def\labelenumi{(\alph{enumi})}
\item
  证明合成 \(xx^{-1}, x^{-1}e_x, e_x'x^{-1}, e_x e_x, e_x' e_x'\) 都有定义，并且 \(xx^{-1} = e_x, x^{-1}e_x = e_x'x^{-1} = x^{-1}, e_x e_x = e_x, e_x' e_x' = e_x'\) 。
\item
  每个幂等元 \(e\) （属于 \(\mathbf{E}\) ，§ 1，第 4 号）是所有 \(x\) （使 \(ex\) 有定义者）的左单位元，并且是所有 \(y\) （使 \(ye\) 有定义者）的右单位元。
\item
  为使合成 \(xy\) 有定义，必要且充分的是 \(x\) 的右单位元与 \(y\) 的左单位元相同（要看出该条件是充分的，利用关系式 \(e_y = yy^{-1}\) ）；若 \(xy = z\) ，则 \(x^{-1}z = y\) , \(zy^{-1} = x\) , \(y^{-1}x^{-1} = z^{-1}\) , \(z^{-1}x = y^{-1}\) , \(yz^{-1} = x^{-1}\) （这些关系式左边出现的合成都是有定义的）。
\item
  对任意两个幂等元 \(e, e'\) （属于 \(\mathbf{E}\) ），设 \(\mathbf{G}_{e, e'}\) 表示这样的 \(x \in \mathbf{E}\) 所成的集合： \(e\) 是左单位元且 \(e'\) 是 \(x\) 的右单位元。证明：若 \(\mathbf{G}_{e, e}\) 非空，则它在由 \(\mathbf{E}\) 上的律诱导的律下是一个群。
\end{enumerate}

E若还满足以下条件，则称为群胚：

\begin{enumerate}
\def\labelenumi{(\arabic{enumi})}
\setcounter{enumi}{3}
\tightlist
\item
  对于所有幂等元 \(e, e', G_{e, e'}\) 是非空的。
\end{enumerate}

\begin{enumerate}
\def\labelenumi{(\alph{enumi})}
\setcounter{enumi}{4}
\item
  在群胚 E 中，若 \(a \in G_{e,e'}\) ，证明 \(x \mapsto xa\) 是 \(G_{e,e}\) 到 \(G_{e,e'}, y \mapsto ay\) 上的一个双射映射，是 \(G_{e',e'}\) 到 \(G_{e,e'}\) 上的一个双射映射，而 \(x \mapsto a^{-1}xa\) 是群 \(G_{e,e}\) 到群 \(G_{e',e'}\) 上的一个同构。
\item
  证明 § 1，习题 4(b) 中定义的律在集合 \(\Psi\) 上确定一个群胚结构，如果族 \(\mathfrak{F}\) 的所有集合都是等势的。
\end{enumerate}

\begin{enumerate}
\def\labelenumi{\arabic{enumi}.}
\setcounter{enumi}{23}
\tightlist
\item
  \begin{enumerate}
  \def\labelenumii{(\alph{enumii})}
  \tightlist
  \item
  给定任意集合 E，在集合 \(E \times E\) 上考虑一个并非处处有定义、以乘法记号书写的合成律，在该律下， \((x, y)\) 与 \((y', z)\) 的合成仅当 \(y' = y\) 时才有定义，且此时其值为 \((x, z)\) 。证明 \(\mathbf{E} \times \mathbf{E}\) 配备该合成律是一个群胚（习题 23）。
  \end{enumerate}
\end{enumerate}

\begin{enumerate}
\def\labelenumi{(\alph{enumi})}
\setcounter{enumi}{1}
\tightlist
\item
  设 \((x, y) \equiv (x', y')\) (R) 是一个与上述合成相容的等价关系（即 \((x, y) \equiv (x', y')\) 与 \((y, z) \equiv (y', z')\) 蕴含 \((x, z) \equiv (x', z'))\) ；进一步假设关系 R 满足以下条件：对所有 \(x, y, z\) ，存在且仅存在一个 \(t \in \mathbf{E}\) 使得 \((x, y) \equiv (z, t)\) ，且存在一个 \(u \in \mathbf{E}\) 使得 \((x, y) \equiv (u, z)\) 。
\end{enumerate}

证明在这些条件下， \(E \times E\) 上的群胚结构关于 R 的商结构是一个群结构（先证明商律处处有定义，因为若 \(\dot{x}, \dot{y}, \dot{z}\) 是三个类，满足 \(\dot{x}\dot{y} = \dot{x}\dot{z}\) ，则 \(\dot{y} = \dot{z}\) ；最后建立：对所有 \(x \in E, y \in E, (x, x) \equiv (y, y)\) ）。

\begin{enumerate}
\def\labelenumi{(\alph{enumi})}
\setcounter{enumi}{2}
\tightlist
\item
  设 G 是通过把上述结构赋予 \(E \times E\) 关于 R 的商而得到的群。设 a 为 E 的任意元素；若对所有 \(x \in E\) ， \(f_{a}(x)\) 表示 \((a, x)\) 的模 R 的类，证明 \(f_{a}\) 是 E 到 G 上的一个双射映射，且关系 \((x, y) \equiv (x', y')\) 等价于关系
\end{enumerate}

\[
f _ {a} (x) \left(f _ {a} \left(x ^ {\prime}\right)\right) ^ {- 1} = f _ {a} (y) \left(f _ {a} \left(y ^ {\prime}\right)\right) ^ {- 1}.
\]

要使 G 是交换群，必要且充分的是关系 \((x,y)\equiv (x',y')\) 蕴含 \((x,x')\equiv (y,y')\) 。

¶ 25. 设 E 是一个集合，f 是 \(E^{m}\) 到 E 的映射；我们记

\[
f (x _ {1}, x _ {2}, \dots , x _ {m}) = x _ {1} x _ {2} \dots x _ {m};
\]

假设 \(f\) 满足以下条件：

\[
(x _ {1} x _ {2} \dots x _ {m}) x _ {m + 1} \dots x _ {2 m - 1} = x _ {1} (x _ {2} x _ {3} \dots x _ {m + 1}) x _ {m + 2} \dots x _ {2 m - 1}\tag{1}
\]

恒成立；

\begin{enumerate}
\def\labelenumi{(\arabic{enumi})}
\setcounter{enumi}{1}
\tightlist
\item
  对于所有 \(a_1, a_2, \ldots, a_{m-1}\) ，映射
\end{enumerate}

\[
x \mapsto x a _ {1} a _ {2} \dots a _ {m - 1}
\]

\[
x \mapsto a _ {1} a _ {2} \dots a _ {m - 1} x
\]

是 E 到 E 上的双射。

\begin{enumerate}
\def\labelenumi{(\alph{enumi})}
\tightlist
\item
  证明恒有
\end{enumerate}

\[
(x _ {1} x _ {2} \dots x _ {m}) x _ {m + 1} \dots x _ {2 m - 1} = x _ {1} x _ {2} \dots x _ {i} (x _ {i + 1} \dots x _ {i + m}) x _ {i + m + 1} \dots x _ {2 m - 1}
\]

对每个指标 \(i\) （ \(1 \leqslant i \leqslant m - 1\) ）成立（对 \(i\) 作归纳论证，考虑元素

\[
\left(\left(x _ {1} x _ {2} \dots x _ {m}\right) x _ {m + 1} \dots x _ {2 m - 1}\right) a _ {1} a _ {2} \dots a _ {m - 1}).
\]

\begin{enumerate}
\def\labelenumi{(\alph{enumi})}
\setcounter{enumi}{1}
\tightlist
\item
  对于所有 \(a_1, a_2, \ldots, a_{m-2}\) ，存在 \(u \in E\) 使得对所有 \(x\)
\end{enumerate}

\[
x = x a _ {1} a _ {2} \dots a _ {m - 2} u = u a _ {1} a _ {2} \dots a _ {m - 2} x.
\]

\begin{enumerate}
\def\labelenumi{(\alph{enumi})}
\setcounter{enumi}{2}
\tightlist
\item
  在集合 \(\mathbf{E}^k\) ——它由诸序列 \((u_1, u_2, \ldots, u_k)\) （ \(k\) 个元素，属于 \(\mathbf{E}\) ）组成——（ \(1 \leqslant k \leqslant m - 1\) ）中，考虑表述如下的等价关系 \(\mathbf{R}_k\) ：对所有 \(x_1, x_2, \ldots, x_{m-k}, u_1 u_2 \ldots u_k x_1 x_2 \ldots x_{m-k} = v_1 v_2 \ldots v_k x_1 x_2 \ldots x_{m-k}\) ；令 \(\mathbf{E}_k\) 表示商集 \(\mathbf{E}^k / \mathbf{R}_k\) ，令 \(\mathbf{G}\) 表示 \(\mathbf{E}_1 = \mathbf{E}, \mathbf{E}_2, \ldots, \mathbf{E}_{m-1}\) 的和集（集合论，II，§4，第 8 号）。设 \(\alpha \in \mathbf{E}_i\) , \(\beta \in \mathbf{E}_j\) ；若 \((u_1, u_2, \ldots, u_i)\) 是类 \(\alpha\) 的一个序列，而 \((v_1, \ldots, v_j)\) 是类 \(\beta\) 的一个序列，考虑序列
\end{enumerate}

\[
\left(u _ {1}, u _ {2}, \dots , u _ {i}, v _ {1}, \dots , v _ {j}\right)
\]

（在 \(\mathbf{E}^{i + j}\) 中，若 \(i + j < m\) ），或者序列

\[
\left(u _ {1}, u _ {2}, \dots , u _ {i + j - m}, \left(u _ {i + j - m + 1} \dots u _ {i} v _ {1} v _ {2} \dots v _ {j}\right)\right)
\]

（在 \(E^{i+j-m+1}\) 中，若 \(i+j\geqslant m\) ）；证明该序列在 \(E_{i+j}\) （相应地 \(E_{i+j-m+1}\) ）中的类仅依赖于类 \(\alpha\) 与 \(\beta\) ；把它记作 \(\alpha.\beta\) 。证明这样就在 G 上定义了一个群律， \(H=E_{m-1}\) 是 G 的一个正规子群，且商群 G/H 是 m-1 阶的循环群；最后证明 E 与生成 G/H 的一个类（模 H ）相同，并且 \(x_{1}x_{2}\ldots x_{m}\) 正是序列 \((x_{1},x_{2},\ldots,x_{m})\) 在群 G 中的合成。

¶ 26. 设 G, G' 是两个群， \(f: G \to G'\) 是一个映射，使得对 G 的任意两个元素 \(x, y\) ， \(f(xy) = f(x)f(y)\) 或 \(f(xy) = f(y)f(x)\) 。我们打算证明： \(f\) 是 G 到 G' 中的一个同态，或是 G 到反群 G'0 中的一个同态（换言之，或者 \(f(xy) = f(x)f(y)\) 对每个有序对 \((x, y)\) 成立，或者 \(f(xy) = f(y)f(x)\) 对每个有序对 \((x, y)\) 成立）。

\begin{enumerate}
\def\labelenumi{(\alph{enumi})}
\item
  证明集合 \(N = f(e')\) （ \(e'\) 是 \(G'\) 的单位元）是 G 的一个正规子群，并且 f 通过 \(G \xrightarrow{p} G/N \xrightarrow{g} G'\) 分解，其中 g 是单射。因此可以只考虑 f 为单射的情形，下文中将始终作此假定。
\item
  证明：如果 \(xy = yx\) 那么 \(f(x)f(y) = f(y)f(x)\) （考虑 \(f(x^2y)\) 和 \(f(x^2y^2)\) ，以多种方式表达它们）。
\item
  证明：如果 \(f(xy) = f(x)f(y)\) 那么 \(f(yx) = f(y)f(x)\) 。（借助 (b) 证明，可以只考虑 \(xy \neq yx\) 的情形。）
\item
  证明 \(f(xyx) = \hat{f}(x)f(y)f(x)\) 。（依 \(xy = yx\) 或 \(xy \neq yx\) 而使用 (a) 或 (b)；在第二种情形，考虑 \(f(x^2y)\) 与 \(f(yx^2)\) 。）
\item
  设 A 为这样的 \(x \in G\) 所成的集合：\(f(xy) = f(x)f(y)\) 对所有 \(y \in G\) 成立；设 B 为这样的 \(x \in G\) 所成的集合：\(f(xy) = f(y)f(x)\) 对所有 \(y \in G\) 成立。证明 \(A \neq G\) 、 \(B \neq G\) 与 \(A \cup B = G\) 是不可能的情形。（根据 (b)，此时将存在 A 中的 \(x, y\) ，使得 \(f(xy) = f(x)f(y) \neq f(y)f(x)\) ，以及 B 中的 \(u, v\) ，使得 \(f(uv) = f(v)f(u) \neq f(u)f(v)\) 。依次考虑 \(xu \in A\) 、 \(xu \in B\) 两种可能性，由此推出矛盾；第一种情形考虑 \(f(xuv)\) ，第二种情形考虑 \(f(yxu)\) 。）
\end{enumerate}

证明的其余部分在于用归谬法证明 \(\mathbf{A} \cup \mathbf{B} \neq \mathbf{G}\) 是不可能的。若 \(\mathbf{A} \cup \mathbf{B} \neq \mathbf{G}\) ，则存在 \(a, b, c\) （属于 \(\mathbf{G}\) ），使得

\[
f (a b) = f (a) f (b) \neq f (b) f (a) \quad \text { and } \quad f (a c) = f (c) f (a) \neq f (a) f (c).
\]

\begin{enumerate}
\def\labelenumi{(\alph{enumi})}
\setcounter{enumi}{5}
\item
  证明 \(f(c)f(a)f(b) = f(b)f(a)f(c)\) 。（依 \(bc \neq cb\) 或 \(bc = cb\) 分两种情形考虑。第一种情形考虑 \(f(bac)\) ；第二种情形利用 (b) 与 (c)，并考虑 \(f(abc), f(bca)\) 与 \(f(bac)\) 。）
\item
  考虑 \(f(abac)\) 并得出矛盾。
\end{enumerate}

\BourbakiExercisesSection

\begin{enumerate}
\def\labelenumi{\arabic{enumi}.}
\tightlist
\item
  在有限群 G 中，证明一个元素 \(a \in G\) 的共轭元素（第 4 号）的个数等于 a 的正规化子的指数，从而是 G 的阶的一个因子。
\end{enumerate}

¶ 2. *若 G 是一个阶为 n 的有限群，则 G 的自同构个数为 \(\leqslant n^{\log n/\log 2}\) （证明存在 G 的一个生成集 \(\{a_{1}, a_{2}, \ldots, a_{m}\}\) ，使得 \(a_{i}\) 不属于由 \(a_{1}, a_{2}, \ldots, a_{i-1}\) 生成的子群（对 \(2 \leqslant i \leqslant m\) ）；由此推出 \(2^{m} \leqslant n\) ，且 G 的自同构个数为 \(\leqslant n^{m})\cdot*\)

\begin{enumerate}
\def\labelenumi{\arabic{enumi}.}
\setcounter{enumi}{2}
\tightlist
\item
  设 \(\Gamma\) 是一个群 \(G\) 的自同构群， \(\Delta\) 是 \(G\) 的内自同构群；证明 \(\Delta\) 是 \(\Gamma\) 的一个正规子群。为使一个自同构 \(\sigma\) （ \(G\) 的）与 \(G\) 的所有内自同构可交换，必要且充分的是：对所有 \(x \in G\) ， \(x^{-1}\sigma(x)\) 属于 \(G\) 的中心；由此推出，若 \(G\) 的中心约化为单位元，则 \(\Delta\) 在 \(\Gamma\) 中的中心化子也约化为单位元。
\end{enumerate}

¶ 4. (a) 设 G 是一个非交换单群，Γ 是 G 的自同构群，Δ 是 G 的内自同构群（同构于 G）。若 s 是群 Γ 的一个自同构，证明 \(s(\Delta) = \Delta\) （利用上面的习题 3 以及 § 4，第 9 号，命题 15，注意交 \(\Delta \cap s(\Delta)\) 不可能仅由 Γ 的单位元组成）。

\begin{enumerate}
\def\labelenumi{(\alph{enumi})}
\setcounter{enumi}{1}
\item
  证明 \(\Gamma\) 的使 \(\Delta\) 的每个元素都不变的唯一自同构是恒等映射（利用该自同构使 \(\alpha\) 与 \(\sigma \alpha \sigma^{-1}\) （对所有 \(\alpha \in \Delta\) 与 \(\sigma \in \Gamma\) ）都不变这一事实，并使用习题 3）。
\item
  设 \(s\) 是 \(\Gamma\) 的一个自同构， \(\phi\) 是同构 \(x \mapsto \operatorname{Int}(x)\) （G 到 \(\Delta\) 上的）， \(\psi\) 是其逆同构， \(\sigma\) 是 G 的自同构 \(\psi \circ s \circ \phi\) ；证明自同构 \(\xi \mapsto \sigma^{-1}s(\xi)\sigma\) （ \(\Gamma\) 上的）是恒等自同构（利用 (b)，并注意到，对所有 \(x \in G\) ， \(s(\operatorname{Int}(x)) = \operatorname{Int}(\sigma(x)))\) 。由此推出， \(\Gamma\) 的每个自同构都是内自同构。
\end{enumerate}

5. 设 G 是一个群。

\begin{enumerate}
\def\labelenumi{(\alph{enumi})}
\item
  若 H 是 G 的一个指数为有限数 n 的子群，证明 H 的诸共轭子群的交 N 的指数是 \(n!\) 的一个约数（注意 G/N 同构于 \(S_{n}\) 的一个子群）。
\item
  G 称为剩余有限的，如果它的有限指数子群的交是 \(\{e\}\) 。证明这等价于说 G 同构于有限群的一个乘积的一个子群。
\end{enumerate}
